\chapter{外设与队列：SSD、硬盘、网卡}

主存、总线和 DMA 已经讲完，本章看设备本身。磁盘按块搬运，网络按帧到达，它们的速度都比内存慢几个数量级，而且慢的方式完全不同：闪存能读快写慢还有寿命问题，机械盘怕寻道，网络怕突发。但它们的接入方式惊人地一致——操作系统准备好队列和描述符，设备通过 DMA 搬运，完成后再通知。理解外设，一半在理解介质，一半在理解这套队列机制怎样把慢介质藏在高吞吐后面。

\begin{longtable}{L{0.16\linewidth}L{0.36\linewidth}L{0.38\linewidth}}
\toprule
设备 & 介质的脾气 & 队列机制要解决什么 \\
\midrule
SSD & 读快写慢、擦除单位大、有寿命 & 写放大、FTL、磨损均衡 \\
机械硬盘 & 寻道和旋转是毫秒级 & 调度合并、NCQ、顺序化 \\
网卡 & 帧突发、线速不可退 & 描述符环、中断聚合、多队列 \\
\bottomrule
\end{longtable}

\begin{keyidea}
外设性能的本质是"让介质按其擅长的方式工作，让队列吸收其不擅长的部分"。读外设文档时，先问介质的物理限制是什么，再看队列机制怎样弥补。
\end{keyidea}

\section{一、块设备与文件系统的硬件接口}

本节回答一个朴素的问题：一次写文件，到了设备那端究竟发生
了什么。答案分两半。一半是编址：块设备从不按字节寻址，
主机看到的是一组从 0 起编的扇区，这套编址从 CHS 几何参
数一路演化到 LBA；另一半是缓存：从 CPU 到介质隔着页缓
存、设备缓存两层 RAM，"写完了"在哪一层被宣布，直接决定
断电后数据还在不在。

\topic{块设备与字节设备的分界}
操作系统把外设粗分成两类。字节设备（character device）
按流读写：串口每来一个中断交出一个字节，键盘把按键编码
逐个上交，终端把输出字节按顺序显示。流没有"位置"可言：
读过的字节不会再来，写出去的字节收不回，"回到第 37 个
字节"这种请求在流上不成立。块设备（block device）按固
定大小的整块读写：磁盘、SSD、eMMC 都以块为最小传输单位，
每块有编号，可按号随机读、随机写、反复改写。前者像自来
水管，后者像一墙编了号的储物柜；\code{/dev} 里设备文件
类型位的 \code{b} 与 \code{c}，标的就是这个分界。

分界不是软件随意划的，它几乎由介质单方面决定。存储介质
普遍成块，因为定位与寻址的代价按次数收、不按字节收：机
械盘的磁头移一次、闪存的页选通一次、磁带走带定位一次，
都是固定开销，摊到越多字节上越划算；介质于是把字节捆成
块、给块编号，让一次定位服务一整块。终端和串口相反：它
们的"介质"是线路另一端的人或进程，数据以事件节奏零散到
达，没有可寻址的静止形态，天然成流。设备分类表因此大体
是一张介质决定论的清单：

\begin{longtable}{L{0.13\linewidth}L{0.25\linewidth}L{0.23\linewidth}L{0.27\linewidth}}
\toprule
类别 & 代表设备 & 读写单位 & 为什么是这样 \\
\midrule
块设备 & 机械盘、SSD、eMMC、U 盘 & 整块，按号寻址，可改写 & 介质按块组织，定位代价按次收 \\
字节设备 & 串口、键盘、终端、打印机 & 字节流，无地址，不可回看 & 数据以事件节奏到达，无静止形态 \\
网络设备 & 网卡 & 帧（分组），突发到达 & 既非流也非块，第四节专门讲 \\
特例 & 内存条 & 按地址随机访问 & 它就是"地址"本身，不需块这层抽象 \\
\bottomrule
\end{longtable}

两处边界值得留神。其一，两种接口可以互相伪装，只是代价
不对等：块设备装成流很容易，\code{dd} 按字节计数拷盘，
靠的是软件把流切成块、逐块下发；字节设备装成块则要在流
上模拟随机访问，理论上可行，工程上没人做。磁带是骑在边
界上的活例子：它按记录块寻址、算块设备，却只支持顺序定
位，跳读远处的块以十秒计，于是当不了文件系统的载体，只
配做归档——"可随机定位"才是块设备配享文件系统的入场券。
其二，"可改写"对块设备是逻辑承诺、不是物理事实：主机以
为自己在原地改写了 37 号块，SSD 的 FTL 完全可能把新数
据写到别处、只改一张映射表。这是下一节的主题，这里先记
住结论：块设备承诺的是"按号寻址的整块读写"，不承诺任何
物理位置。

还要防止一个误读：流设备也有缓冲，缓冲不改变流的本质。
串口的 UART 自带 16 字节 FIFO，终端有行缓冲，管道在内核
里有 64 KiB 环形缓冲——它们吸收的是速率抖动，不是引入
地址：缓冲里的字节依然先到先走、读完即弃，没有任何"回
到第 37 字节"的办法。判断一个设备是流还是块，不问它有
没有缓冲，只问它的数据有没有编号、能不能按编号回头取。

\topic{扇区、块与 LBA}
块设备的最小寻址单位叫扇区（sector），称呼来自机械盘：
盘片上每个磁道被切成若干弧段，一段就是一个扇区。早期接
口把物理几何直接暴露给软件——第六章 int 13h 的调用约定
就是标本：柱面号、磁头号、扇区号分装三个寄存器，柱面
10 位、磁头 8 位、扇区 6 位且从 1 起编。这套 CHS
（cylinder-head-sector）编址把容量上限锁死在寄存器位宽
里：

\begin{lstlisting}
CHS 上限 = 1024 柱面 x 256 磁头 x 63 扇区 x 512 字节
         = 8,455,716,864 字节，约 8.4 GB
DOS 兼容约定把磁头数限到 255，上限收缩到约 7.8 GiB
\end{lstlisting}

更早还有一道更矮的墙：BIOS 与 IDE 寄存器各自的几何位宽
取交集，最窄处拼出的上限更小：

\begin{longtable}{L{0.26\linewidth}L{0.16\linewidth}L{0.16\linewidth}L{0.2\linewidth}}
\toprule
接口 & 柱面位数 & 磁头位数 & 扇区位数 \\
\midrule
BIOS int 13h & 10 & 8 & 6（从 1 起编） \\
IDE/ATA & 16 & 4 & 8 \\
两者交集（实际可用） & 10 & 4 & 6 \\
\bottomrule
\end{longtable}

交集那行就是 $1024\times16\times63\times512$ 字节
$\approx504$ MiB 这道墙的出处，1990 年代中期的盘纷纷撞
线。解法不是继续加位宽，而是换掉抽象本身：LBA（logical
block addressing）把盘看成从 0 起编的一维扇区数组，几何
换算藏进盘内控制器。int 13h 扩展（\code{AH=0x42}）改用
磁盘地址包传 64 位 LBA；ATA 协议的 LBA28 给 28 位地址，
上限 $2^{28}\times512$ 字节 $=128$ GiB；LBA48 给 48 位，
上限 $2^{57}$ 字节 $=128$ PiB，容量焦虑就此退场。CHS 从
此退化为兼容字段：现代盘报出的柱面磁头数是编出来哄老软
件的，与盘片真实布局毫无关系。第六章讲固件把硬件细节封
装在寄存器约定后面，LBA 是同一招用在编址上——把"盘面
几何"换成"线性数组"，接口从此与介质解耦。

扇区本身多大，三十年里的答案是 512 字节；2010 年前后
Advanced Format 把物理扇区扩到 4 KiB：同样的盘片面积能
放下更多数据（每扇区的同步头、间隙与纠错码占比随扇区变
大而变小），纠错能力也更强。麻烦在操作系统生态仍按 512
字节发命令，过渡期于是生出两种盘：

\begin{longtable}{L{0.1\linewidth}L{0.24\linewidth}L{0.52\linewidth}}
\toprule
盘类型 & 逻辑/物理扇区 & 部分写的代价 \\
\midrule
512n & 512 B / 512 B & 无翻译，老盘直接写 \\
512e & 512 B / 4 KiB & 固件翻译；写不满一个物理扇区时盘内读-改-写 \\
4Kn & 4 KiB / 4 KiB & 无翻译；只认 512 B 的老软件直接出错 \\
\bottomrule
\end{longtable}

512e 引出了"原子单位"这个关键概念。块设备的读写以逻辑
扇区为原子单位：一次写要么整个扇区成功、要么整个扇区失
败，不存在"半个新扇区加半个旧扇区"的可见中间态（断电恰
好发生在两个扇区之间是另一回事，靠文件系统日志兜底）。
主机写一个 512 字节逻辑扇区、物理扇区却是 4 KiB 时，那
4 KiB 里有八分之七是必须保留的老数据，盘内只能读-改-
写：先把物理扇区整页读进盘内缓存，替换其中 512 字节，
再整页写回。一次逻辑写放大成"一次读加一次写"；在
7200 rpm 的机械盘上若错过当圈的写窗口，还要多等一整圈
$60/7200\approx8.3$ ms。读-改-写（read-modify-write）
这个词后面还会反复出现：文件系统未对齐写、RAID 的部分
条带写、闪存的页内改写，全是同一机制在不同层级的化身。

还要记住 LBA 的编号单位是逻辑扇区、不是字节：同一个
LBA 2048，在 512e 盘上是字节偏移 1 MiB，在 4Kn 盘上是
8 MiB。分区表、引导代码、\code{dd} 的 \code{skip} 计数，
全都按逻辑扇区计——换盘或换镜像格式时，这是最容易算错
的一位。

\topic{文件系统块与设备扇区}
文件系统在块设备之上又定了一级自己的单位：文件系统块，
Linux 桌面与服务器最常见的取值是 4 KiB。文件的分配、空
闲位图、日志全按这个粒度登记，一次写文件最终化成若干个
4 KiB 块下发到块层。两层单位同名不同级：设备扇区是硬件
的原子单位，文件系统块是软件的分配单位。两者对齐则相安
无事，错开则每次写都在交学费。

对齐是道算术题。设物理扇区 4 KiB，分区从第 63 个 512 字
节扇区开始——MBR 时代按"柱面对齐"留下的经典起点：
$63\times512=32\,256$ 字节，除以 4096 得 7.875，不是整
数。于是分区里每个 4 KiB 文件系统块都横跨两个物理扇区：
写这样一个块要动两个物理扇区，在 512e 盘上每个被部分覆
盖的物理扇区都要读-改-写；一次文件系统写拆成两笔、每
笔先读后写，随机写吞吐近乎腰斩，顺序写也因多出的盘内操
作而打折。闪存上还有第二重错位：擦除块以几百 KiB 计
（下一节细讲），分区起点与擦除块边界错开，垃圾回收每轮
都要多搬一批页。

解决办法笨拙而有效：对齐到一个足够大的公倍数。1 MiB
$=2^{20}$ 字节，是 512 字节扇区的 2048 倍、4 KiB 扇区的
256 倍，也能整除常见 RAID 条带（64/128/256 KiB）和多数
闪存擦除块（256 KiB-1 MiB）。2000 年代末起，主流分区工
具把默认起点从 63 号扇区改到 2048 号扇区（整 1 MiB），
"分区对齐 1 MiB"从此成为惯例。三种起点的对照一目了然：

\begin{longtable}{L{0.24\linewidth}L{0.16\linewidth}L{0.2\linewidth}L{0.26\linewidth}}
\toprule
分区起点 & 起点字节 & 整除 4 KiB？ & 后果 \\
\midrule
63 号 512 B 扇区 & 32 256 B & 否（7.875） & 每个文件系统块横跨两个物理扇区 \\
64 号 512 B 扇区 & 32 768 B & 能（8） & 对齐扇区，但不对齐 1 MiB \\
2048 号 512 B 扇区 & 1 MiB & 能（256） & 扇区、条带、擦除块全部对齐 \\
\bottomrule
\end{longtable}

检查办法只剩一句算术：分区起始 LBA 乘扇区大小，能同时被
4096、条带大小与擦除块大小整除，才算对齐；\code{fdisk -l}
与 \code{lsblk} 都能直接读出起点，装完系统第一件事就该
验这道题。

\topic{读写缓存的层级}
从 CPU 到介质，数据要经过两层 RAM，每一层都可能"代收"
一次写：

\begin{longtable}{L{0.14\linewidth}L{0.26\linewidth}L{0.2\linewidth}L{0.24\linewidth}}
\toprule
层级 & 位置与容量级 & 读命中时延量级 & 断电后 \\
\midrule
页缓存 & 主存，GiB 级 & 约 100 ns & 丢失 \\
设备缓存 & 盘内 DRAM，MiB-GiB 级 & 数十 $\mu$s & 无掉电保护则丢失 \\
介质 & 闪存或盘片，百 GiB-TiB 级 & $10^{2}$ $\mu$s-$10$ ms & 保留 \\
\bottomrule
\end{longtable}

页缓存（page cache）是操作系统拿空闲主存开的代收点：写
系统调用把数据拷进页缓存即返回，真正的写盘由回写线程按
几十秒级的节奏异步完成。第五章讲的 DMA 与描述符环就在
回写这一刻登场：内核把脏页排进块层队列、写成描述符，磁
盘控制器按描述符用 DMA 搬数据，搬完写回完成状态，CPU
全程不经手字节。读的路径对称：命中页缓存的读是百纳秒
级，设备连被惊动都没有；未命中才排进队列、走一遍 DMA，
操作系统还会顺手预读（readahead），把顺序后继的块一并
读进页缓存。这台机器上"读文件"与"读盘"是两回事，日常
十有八九是前者。三个时延层级 $10^{2}$ ns、$10^{2}$
$\mu$s、$10$ ms 排成 $1:10^{3}:10^{5}$ 的阶梯，正是第五
章存储层次向外设方向的延伸。

设备缓存是盘内的一小块 DRAM（机械盘几十到几百 MiB，SSD
可达 GiB 级），职责有二：读预取，把顺序后继扇区提前读进
来；写暂存，把主机刚写来的数据先存下。SSD 的盘内缓存还
有一份差事：下一节讲的 FTL 映射表就放在里面。写暂存的
策略由 ATA 的 WCE（write cache enable）一类的开关控制。
开，对应 write-back 语义：数据一进盘内缓存，控制器就向
上报完成，随后择机写介质——报完成那一刻数据还不在介质
上，断电即丢，除非盘内有掉电保护。关，对应 write-through
语义：报完成之前数据必须已在介质上，每笔写慢一个介质写
时延，换来"完成即安全"。主机侧有同名概念：页缓存本身就
是操作系统的 write-back，\code{O\_DIRECT} 与
\code{O\_SYNC} 是绕过或收紧它的两个阀门。于是"写命中哪
一层才算安全"有了明确答案：数据到达非易失的去处——介
质本身，或带掉电保护、保证能把缓存清空的盘内 RAM——这
次写才算数；停在任何一层易失 RAM 上的"完成"，都只是承
诺。

这个答案解释了企业盘与消费盘的一处规格差异。企业级 SSD
板上留一组储能电容——第八章讲过去耦电容按时间尺度分层，
这组电容接管其中最慢的一档：市电消失——保证突然断电时
盘内缓存能全部写进闪存。有了它，盘可以放心用 write-back
报完成，既快又安全；没有它的消费盘，write-back 的"完成"
在断电瞬间可能兑不了现。

两层缓存的性格也不同。页缓存面向工作集：热文件反复命中，
命中与否取决于访问的局部性；设备缓存面向流：机械盘的预
取对顺序读几乎百发百中，对随机读基本帮不上忙，于是盘
"顺序快、随机慢"的脾气在缓存层又被放大了一次。

\topic{sync/fsync 的硬件含义}
页缓存异步回写意味着 \code{write} 返回时，数据多半还躺
在主存里。要让数据真正到达介质，必须有人把每一层都走完，
这正是 \code{fsync} 的语义。一次 \code{fsync(fd)} 做三件
事：把该文件在页缓存里的脏页排进块层队列，等它们全部写
完；向设备发 FLUSH CACHE 命令，要求盘把缓存里所有暂存
写清到介质；等设备确认 FLUSH 完成，系统调用才返回。
\code{sync} 是同一件事的全盘版：清所有脏页、刷所有设备
的缓存。FLUSH 在不同协议里名字不同，角色是同一个：

\begin{longtable}{L{0.16\linewidth}L{0.34\linewidth}L{0.36\linewidth}}
\toprule
协议 & 命令名 & 语义 \\
\midrule
ATA/SATA & FLUSH CACHE（EXT） & 把盘内写缓存全部写入介质 \\
SCSI/SAS & SYNCHRONIZE CACHE & 同上，可指定 LBA 区间 \\
NVMe & Flush & 把指定命名空间的易失缓存写清 \\
\bottomrule
\end{longtable}

注意 FLUSH 是发给设备的命令，走设备自己的队列：在第五章
的描述符环里，它与读写命令一样占一个描述符，只是语义从
"搬数据"变成"把已搬的数据做实"。

fsync 有多贵，拆开算。极端情形先看 7200 rpm 机械盘：脏
页在盘上随机分布，每一笔都要定位——平均寻道约 4 ms，平
均旋转等待半圈 $60/7200/2\approx4.17$ ms，4 KiB 传输约
27 $\mu$s——单笔约 8-12 ms，FLUSH 再等最后一笔真正写
入磁性介质。单线程每提交一个事务就 fsync 一次，吞吐上限
被钉在 $1/10\ \mathrm{ms}=100$ 次/秒量级。换无掉电保护
的 SATA SSD：系统调用与文件系统日志处理约 10 $\mu$s，命
令下发加 DMA 搬运 4 KiB 约 10 $\mu$s，FLUSH 要等一次
TLC 页编程、约 1 ms（为什么这么慢，下一节讲），合计仍是
毫秒级，上限约 $10^{3}$ 次/秒。带掉电保护的盘可以把
FLUSH 截停在缓存层——电容兜底，缓存就算数——fsync 收
缩到几十微秒，上限抬高两个量级。一次 fsync 的时间构成
列表如下：

\begin{lstlisting}
fsync 时间构成（无掉电保护的 SATA SSD，量级）
  系统调用 + 文件系统日志写页缓存   约 10 us
  命令下发 + DMA 搬运 4 KiB        约 10 us
  FLUSH：等 TLC 页编程写入介质     约 1 ms
  合计                             毫秒级，上限约 10^3 次/秒
\end{lstlisting}

把这条层级链与 fsync 的走法画在一起，普通写与同步写
的分岔点一目了然：

\begin{pdfdiagram}
  \node[hw state, minimum width=4.6cm] (app) at (2.6,0) {应用：\code{write} / \code{fsync}};
  \node[hw memory, minimum width=4.6cm] (pc) at (2.6,-1.5) {页缓存：主存 GiB 级，读命中约 100 ns};
  \node[hw memory, minimum width=4.6cm] (dc) at (2.6,-3.0) {设备缓存：盘内 DRAM，数十 $\mu$s};
  \node[hw memory, minimum width=4.6cm] (med) at (2.6,-4.5) {介质：闪存或盘片，$10^{2}$ $\mu$s--$10$ ms};
  \draw[hw bus] (app.south) -- (pc.north);
  \draw[hw bus] (pc.south) -- (dc.north);
  \draw[hw bus] (dc.south) -- (med.north);
  \node[hw label, anchor=east] at (2.35,-0.75) {\code{write} 拷进即返回};
  \node[hw label, anchor=east] at (2.35,-2.25) {回写线程：几十秒级 DMA};
  \node[hw label, anchor=east] at (2.35,-3.75) {FLUSH 后写入介质};
  % fsync 穿透路径：右侧虚线一路向下直到介质
  \draw[hw return] (app.east) -- (6.6,0) -- (6.6,-4.5) -- (med.east);
  \node[hw label, anchor=west] at (6.8,-0.75) {① 脏页排进块层队列等写完};
  \node[hw label, anchor=west] at (6.8,-2.25) {② FLUSH 清设备缓存};
  \node[hw label, anchor=west] at (6.8,-3.75) {③ 设备确认后才返回};
  \node[hw label] at (3.3,-5.35) {普通写停在页缓存即算完成；\code{fsync} 一路穿透到介质——停在易失 RAM 的"完成"只是承诺};
\end{pdfdiagram}
\diagramnote{块设备访问的三层结构与 fsync 的穿透路径。写系统调用把数据拷进页缓存即返回，回写线程按几十秒级的节奏异步写盘，盘内还有一层设备缓存代收；fsync 则把每一层都走完：脏页排进块层队列等写完、发 FLUSH 命令清设备缓存、等设备确认才返回。数据到达非易失的去处才算数——停在任何一层易失 RAM 上的"完成"，都只是承诺。}

数据库对 fsync 延迟敏感，原因全在这笔明细里。事务日志
（WAL）的正确性规则是"日志先到介质、数据后到介质"，每
条提交记录都要过 fsync 这一关，提交时延里 fsync 常占九
成以上；它还是串行瓶颈，单个连接的提交速率就是 fsync 速
率的倒数。工程对策全在这笔明细里做文章。组提交把一批事
务攒进同一次 fsync：10 ms 内攒 200 个事务，吞吐就是
$200/10\ \mathrm{ms}=2\times10^{4}$ 事务/秒——单个事务
多等几毫秒，总量放大两百倍，第十章"用时延换吞吐"的交换，
这里是最贵的一例。换带掉电保护的盘，直接砍掉明细里最贵
的那一项。放宽同步频率（如
\code{innodb\_flush\_log\_at\_trx\_commit=2}，每秒刷一次）
则是拿一秒的丢失窗口换两个量级。理解 fsync 的硬件含义之
后，每个调优选项都能换算成"哪一层 RAM 被信任、丢失窗口
有多大"，而不是背参数口诀。

还有一层容易漏算：文件系统自己的日志。ext4 默认
\code{data=ordered} 模式下，\code{fsync} 不只刷目标文件
的数据——若这次写伴随元数据变更（扩块、改大小），还要
先等日志提交写入介质，再刷数据区，一次 \code{fsync} 可
能触发的不止一次 FLUSH。数据库干脆自己管理日志、把数据
文件用 \code{O\_DIRECT} 打开，连页缓存这层也省掉，只留
设备缓存一层要刷——层级少一层，耗时构成就少一项不确定
性。

\section{二、SSD 与闪存的硬件特性}

第一节把块设备当成"按号寻址的线性数组"，那是接口视角。
本节打开 SSD 的盖子看介质：NAND 闪存的读、写、擦三个操
作各有各的粒度，谁也不迁就谁——读按页、写按页、擦按块，
改写不许原地进行。为了让线性数组的承诺在这种介质上成立，
SSD 内部跑着一整套地址翻译、垃圾回收与磨损均衡软件，盘
里那颗控制器的算力超过第七章的整机。

\topic{NAND 闪存的页写块擦}
NAND 闪存的存储单元是浮栅（或电荷俘获）晶体管：绝缘层
里困住的电子改变阈值电压，读出时给控制栅加参考电压、看
沟道导不导通，就知道这一位是 0 还是 1。一个单元存几位，
决定同一片硅的容量、速度与寿命：

\begin{longtable}{L{0.1\linewidth}L{0.14\linewidth}L{0.12\linewidth}L{0.26\linewidth}L{0.24\linewidth}}
\toprule
类型 & 每单元位数 & 阈值档位 & P/E 循环量级 & 典型去处 \\
\midrule
SLC & 1 & 2 & 约 $10^{5}$ 次 & 盘内缓存、企业特殊款 \\
MLC & 2 & 4 & 约 $3\times10^{3}$-$10^{4}$ 次 & 早期主流，现多见于工业盘 \\
TLC & 3 & 8 & 约 $10^{3}$-$3\times10^{3}$ 次 & 当前消费与数据中心主流 \\
QLC & 4 & 16 & 数百次 & 大容量、读多写少 \\
\bottomrule
\end{longtable}

位数越多，阈值电压要分清的档位越多：TLC 是 $2^{3}=8$ 档，
QLC 是 16 档。写入时把电压逐档调准的时间越长，读出时分
辨档位的余量越小，单元能经受的擦写次数也越少——容量、
速度、寿命在同一条曲线上取舍，没有白捡的密度。闪存还有
NOR 一派：支持按字节随机读，用来存固件，但密度低、成本
高；U 盘、SD 卡、eMMC、SSD 全是 NAND 一派——牺牲随机字
节访问换密度，再用控制器把粒度差藏起来。本章说的"闪存"
默认指 NAND。

怪脾气集中在三个操作的粒度上，而且粒度由物理直接给定：

\begin{longtable}{L{0.12\linewidth}L{0.26\linewidth}L{0.2\linewidth}L{0.3\linewidth}}
\toprule
操作 & 粒度 & 时延量级 & 物理原因 \\
\midrule
读 & 页，4-16 KiB & 25-100 $\mu$s & 一页共享字线，整页同时选通 \\
编程（写） & 页，只能 1 写成 0 & 数百 $\mu$s 到约 1 ms & 逐档校准阈值电压，TLC 调 8 档 \\
擦除 & 块，256 KiB-4 MiB & 数 ms & 高压加在整个块的衬底上，按块放电子 \\
\bottomrule
\end{longtable}

编程只能把 1 写成 0，想让任何一位回到 1 必须擦除；擦除
高压施加在整个块共用的衬底上，单独擦一页等于给每页配隔
离与高压通路，芯片面积与可靠性都不答应——块做得越大，
摊到每字节上的外围电路越少，擦除块一路长到数 MiB 正是
这条经济学。三条合起来就是"页写块擦"：写以页为单位，一
页在两次擦除之间只能写一次；擦以块为单位，一块几百页同
生共死。"改写"这个词在闪存里因此失去含义：把一页从 A 改
成 B 不能原地进行，原页里的 0 变不回 1。唯一的出路是把
B 写到另一页、把旧页标记为无效，攒够一整块无效页后一起
擦除。闪存天生是一种"只许追加、不许原地改"的介质。与主
存对照数量级：读一页约 50 $\mu$s，是 DRAM 读的约 500 倍；
写一页约 1 ms，是 DRAM 写的约一万倍——SSD"读快写慢"的
物理出处就在这里。

寿命限制同样来自物理。每次擦除的高压都磨损浮栅外的隧穿
氧化层，磨到电子关不住，块就报废。注意寿命的单位是每块：
一块 2 MiB 的块按 1000 次 P/E 计，一生承担 2 GiB 写入；
整盘寿命是全部块寿命之和。这个"按块计"给后面的磨损均衡
出了题目：谁让擦除集中到少数块上，谁就先杀死整块盘。另
有两条次要脾气记在这里：同一块被读得太多会扰动相邻页
（读干扰），数据放久了电荷会漏（保持力）——FTL 因此要
定期巡检、搬动刷新，这是盘内又一份主机看不见的日常工
作。

\topic{FTL：闪存翻译层}
主机要的接口是"按 LBA 原地改写的线性数组"，介质给的能
力是"只许追加、按块擦除"，中间的鸿沟由 FTL（flash
translation layer，闪存翻译层）填平。FTL 是跑在盘内控制
器上的固件，日常工作四件事，每件对应介质的一条限制：

\begin{longtable}{L{0.16\linewidth}L{0.3\linewidth}L{0.4\linewidth}}
\toprule
职责 & 对应的介质限制 & 做法 \\
\midrule
地址映射 & 不许原地改写 & 写新页、改映射、旧页标脏 \\
空页供给 & 写前必须有已擦空块 & 预留空闲块池，随写随补 \\
垃圾回收 & 擦按块、块内混有有效页 & 搬走有效页，再整块擦除 \\
掉电恢复 & 映射表放在易失 DRAM & 定期写回闪存，凭日志重建 \\
\bottomrule
\end{longtable}

第一件是映射表。FTL 维护从逻辑页号（LBA 按 4 KiB 折算）
到物理页号的映射；主机写某个逻辑页，FTL 从不原地改，而
是把数据写进当前可写块的下一个空页，把映射指向新页、旧
页标脏。"改写 = 写新页 + 旧页标脏"，原地改写就这样被翻
译成追加写。映射的粒度也有讲究：页级映射最灵活、表也最
大；块级映射表小，但块内页偏移被锁死，改一页要动整块，
现代盘几乎清一色选页级。表的体积可以算：256 GiB 的盘按
4 KiB 分页，共 $2^{38}/2^{12}=2^{26}\approx6.7\times10^{7}$
个逻辑页，每条映射 4 字节，整表 $2^{28}$ 字节 = 256 MiB
——恰好是"每 GiB 容量配 1 MiB 表"的比例。这就是 SSD 要
配 DRAM 的首要原因，256 GiB 的盘配 256 MiB 上下；无 DRAM
的廉价盘把映射挤在小块 SRAM 和闪存里，或借主机内存
（NVMe 的 HMB 机制）暂存，随机访问映射本身成了新的瓶颈。

把介质的页块格局与 FTL 的翻译过程画在一起，"改写"在盘内
的真实形态就清楚了：

\begin{pdfdiagram}
  \node[hw state, minimum width=2.0cm] (host) at (1.05,0.35) {主机：改写逻辑页 L};
  \node[hw compute, minimum width=3.4cm] (ftl) at (5.3,0.35) {FTL 映射表\\逻辑页 $\to$ 物理页};
  \draw[hw bus] (host.east) -- (ftl.west);
  \node[hw label, anchor=south] at (3.3,0.9) {改写请求};
  % 块 A：原页所在块，内含若干页
  \draw[BookAmber, line width=0.42pt, rounded corners=2pt] (0.5,-3.2) rectangle (4.1,-1.6);
  \node[hw label] at (2.6,-1.85) {原页所在块};
  \node[hw memory, minimum width=1.0cm, minimum height=0.66cm, fill=BookRed!10] (oldpage) at (1.15,-2.55) {旧页 L};
  \node[hw memory, minimum width=0.8cm, minimum height=0.66cm] at (2.35,-2.55) {页};
  \node[hw memory, minimum width=0.8cm, minimum height=0.66cm] at (3.45,-2.55) {页$\cdots$};
  % 块 B：当前可写块，只许追加
  \draw[BookAmber, line width=0.42pt, rounded corners=2pt] (6.1,-3.2) rectangle (9.7,-1.6);
  \node[hw label] at (8.05,-1.85) {当前可写块};
  \node[hw memory, minimum width=1.0cm, minimum height=0.66cm, fill=BookTeal!15] (newpage) at (6.85,-2.55) {新页 L};
  \node[hw memory, minimum width=0.8cm, minimum height=0.66cm] at (8.05,-2.55) {空页};
  \node[hw memory, minimum width=0.8cm, minimum height=0.66cm] at (9.15,-2.55) {空页$\cdots$};
  % FTL 的两条动作线：共享竖直干线，按 y 分层各自折出
  \draw[hw bus] (ftl.south) -- (5.3,-1.1) -- (6.85,-1.1) -- (newpage.north);
  \node[hw label, above] at (6.9,-1.1) {① 写新页、映射改指};
  \draw[hw return] (ftl.south) -- (5.3,-0.72) -- (1.15,-0.72) -- (oldpage.north);
  \node[hw label, above] at (3.2,-0.72) {② 旧页标脏};
  \node[hw label, text width=9.4cm] at (5.2,-3.85) {块（256 KiB--4 MiB）内含数百页（4--16 KiB）；无效页攒满一整块后，垃圾回收搬走仍有效的页、整块擦除、收回空块。};
\end{pdfdiagram}
\diagramnote{NAND 闪存的页写块擦与 FTL 地址映射。主机以为自己原地改写了逻辑页 L；盘内实际是两步：FTL 把数据写进当前可写块的下一个空页、映射改指新页，旧页标脏等待回收。无效页攒满一整块后，垃圾回收搬走仍有效的页、整块擦除、收回空块——追加写加周期性回收，就是 SSD 在只许追加的介质上提供原地改写接口的全部机制。}

第二件是空页供给：追加写需要源源不断的已擦空块，FTL 平
时预留若干，写完再补。第三件因此必然登场——垃圾回收
（garbage collection，GC）：无效页攒多了，挑一个无效页
占比高的块，把仍有效的页读出、搬进新块，再整块擦除回收。
读、搬、擦全在盘内悄悄进行，主机只看到 LBA 数组一切如常，
只是时延偶尔抖一下。第四件是掉电恢复：映射表在 DRAM 里，
FTL 要定期把表与变更日志写回闪存，断电后凭日志重建，否
则一次意外掉电就是一块砖。

四件事合起来看，SSD 里的控制器名不副实：它远不止"控制"，
那是一颗多核 ARM 级处理器，跑着带映射管理、GC、磨损均衡、
坏块替换、掉电重建的完整软件，配的 DRAM 比第七章整台教
学机的主存大几个数量级。买一块 SSD 等于买回一台专用小计
算机，它的全部工作是让 NAND 假装成一块可以原地改写的磁
盘。下一节会看到，机械盘的控制器同样不闲，只是忙的方向
相反——不是躲开介质的限制，而是迎合介质的转动。

\topic{写放大的计算}
GC 搬运的有效页本身是额外的闪存写入：主机没让它们发生，
FTL 为了让介质继续可用不得不写。写放大（write
amplification，WA）把这笔开销量化：

\begin{lstlisting}
WA = 实际写入闪存的总量 / 主机写入的总量
\end{lstlisting}

WA = 1 是理想值：主机写多少，介质写多少，GC 零搬运。顺
序写满一块再整块改写的负载就接近它——改写让整块的页同
时失效，回收时没有有效页可搬，擦掉即可。WA 大于 1 的部
分，全是 GC 搬家的运费。

具体算一次。设块含 256 页，某块一半（128 页）有效、一半
已无效。回收它：128 个有效页读出并写入新块，闪存写入
128 页；擦掉旧块得 256 个空页，其中 128 页被搬来的数据
占用，剩 128 页供主机写新数据。等主机写满这 128 页盘点：
主机写入 128 页，闪存实际写入 $128+128=256$ 页，WA
$=256/128=2$。有效页占一半，写放大恰为 2 倍。一般化只
剩一步：块内有效页比例为 $v$ 时，回收一块要搬 $vN$ 页、
腾出 $(1-v)N$ 页给主机：

\begin{lstlisting}
WA = (vN + (1-v)N) / ((1-v)N) = 1 / (1 - v)
\end{lstlisting}

$v$ 每涨一点，代价不是线性涨，是倒数式地涨：

\begin{longtable}{L{0.22\linewidth}L{0.14\linewidth}L{0.48\linewidth}}
\toprule
有效页比例 $v$ & WA & 直观画面 \\
\midrule
0 & 1 & 无页可搬，擦掉即可（顺序改写的尽头） \\
0.25 & 约 1.33 & 搬 1 页换 3 页空位 \\
0.5 & 2 & 搬一页换一页，写一倍送一倍 \\
0.9 & 10 & 接近写满的盘做随机写 \\
\bottomrule
\end{longtable}

把 $v=0.5$ 那一行画成过程图，"搬一页换一页"的运费
就看得见了：

\begin{pdfdiagram}
  % 待回收块
  \node[hw label, anchor=south] at (2.1,0.08) {待回收块：256 页};
  \draw[BookAmber, line width=0.42pt, rounded corners=2pt] (0.4,-0.2) rectangle (3.8,-2.0);
  \node[hw memory, minimum width=1.4cm, minimum height=0.72cm, fill=BookTeal!15] (va) at (1.3,-1.3) {有效\\128 页};
  \node[hw memory, minimum width=1.4cm, minimum height=0.72cm, fill=BookRed!10] (ia) at (3.0,-1.3) {无效\\128 页};
  % 新块
  \node[hw label, anchor=south] at (8.9,0.08) {新块};
  \draw[BookAmber, line width=0.42pt, rounded corners=2pt] (7.2,-0.2) rectangle (10.6,-2.0);
  \node[hw memory, minimum width=1.4cm, minimum height=0.72cm, fill=BookTeal!15] (vb) at (8.1,-1.3) {搬入\\128 页};
  \node[hw memory, minimum width=1.4cm, minimum height=0.72cm] (hb) at (9.8,-1.3) {主机新写\\128 页};
  % ① 搬运：左侧出框、顶部通道，再竖直落入新块
  \draw[hw bus] (va.west) -- (0.22,-1.3) -- (0.22,0.72) -- (8.1,0.72) -- (vb.north);
  \node[hw label, anchor=south] at (4.7,0.78) {① 有效页读出、写入新块：闪存写 128 页};
  % ② 主机写满余下空页
  \draw[hw bus] (9.8,-2.75) -- (hb.south);
  \node[hw label, anchor=north] at (9.8,-2.85) {② 主机写入 128 页};
  % ③ 整块擦除
  \draw[hw bus] (2.1,-2.0) -- (2.1,-3.05);
  \node[hw label, anchor=west] at (2.3,-2.55) {③ 整块擦除};
  \draw[BookRule, line width=0.42pt, dashed, rounded corners=2pt] (0.6,-3.15) rectangle (3.6,-3.95);
  \node[hw label] at (2.1,-3.55) {空块：256 页全部可写};
  % 盘点
  \node[hw label] at (5.5,-4.6) {主机写入 128 页，闪存实际写入 $128+128=256$ 页 $\Rightarrow$ WA $=256/128=2$};
\end{pdfdiagram}
\diagramnote{一次 GC 的完整过程（块含 256 页、一半有效）。回收：128 个有效页读出并写入新块，这是主机没有要求的闪存写入；旧块整块擦除，换回 256 个空页，其中 128 页被搬来的数据占用，剩 128 页供主机写新数据。等主机写满盘点：主机写 128 页，介质写 256 页，WA 恰为 2——有效页占一半，写一倍送一倍。}

随机写为什么放大更严重，看 $v$ 的下场就知道。顺序写把失
效集中到少数块上：文件被顺序覆盖，整块一起变脏，$v$ 趋
于 0、WA 趋于 1。均匀随机写把失效撒到所有块上：每个块
都攒一点无效页，又都剩一堆有效页，长期稳态下各块的 $v$
收敛到同一高位——盘越满，$v$ 越高。剩余空间只有一成的
盘，各块平均 $v\approx0.9$，按上式 WA 逼近 10：主机每写
1 GiB，介质实际承受约 10 GiB。GC 挑块因此遵循贪心原则：
优先回收无效页最多的块，同样换回一个空块，$v=0.1$ 的块
只搬一成，$v=0.9$ 的块要搬九成。GC 的时机也影响体验：
空闲时做的后台 GC 主机无感，写路径上被空块耗尽逼出来的
同步 GC 则直接拖慢当次写——同一笔写放大，前者藏在空闲
里，后者变成时延尖峰。同一个"随机写"，在第十章只是时延
问题，在闪存这里同时是寿命问题：放大几倍，P/E 预算就被
多吃几倍。工程上的缓冲垫是预留空间（overprovisioning）：
盘内实装闪存多于标称容量，消费盘多留约 7\%，企业盘留到
28\% 上下，差额直接压低 $v$ 的有效值，把稳态 WA 降回个
位数。消费 TLC 盘还有一招 SLC 缓存：拿一部分闪存按 SLC
模式用，突发写先进缓存、空闲时再腾进 TLC，缓存打穿后的
"缓外速度"才是介质的真速度。

\topic{磨损均衡与寿命}
写放大吃掉的是寿命预算，而 P/E 次数按块计：写入分布不
均，少数块先被擦穿，整盘提前报废。磨损均衡（wear
leveling）的任务是把擦除次数在所有块之间摊平。动态磨损
均衡只盯活跃数据：每次分配空块时挑擦除次数最少的，新写
入自然流向年轻块。它对冷数据无效——一批文件写进去三年
不动，占着的块擦除次数恒为零，剩下的块继续挨打。静态磨
损均衡连冷数据也搬：定期把长期不改写的数据挪到擦除次数
多的块里"养老"，把年轻块换出来轮换。消费盘多做动态加有
限的静态，企业盘两者做足；注意均衡搬运同样计入写放大，
寿命的公平本身是花钱买的。坏块管理走的是同一条备用通道：
出厂坏块与使用中长出的坏块都由备用块池顶替，顶替一次，
可用块与预留空间就少一分。

把擦除次数按块排成直方图，均衡与不均衡的差距一眼见底：

\begin{pdfdiagram}
  % 左：只有动态均衡
  \node[hw label] at (2.55,2.85) {只有动态均衡：热块先被擦穿};
  \draw[BookMuted, line width=0.6pt] (0.4,0) -- (4.7,0);
  \filldraw[fill=BookRed!25, draw=BookRed, line width=0.5pt] (0.475,0) rectangle (0.925,2.5);
  \filldraw[fill=BookTeal!25, draw=BookTeal, line width=0.5pt] (1.095,0) rectangle (1.545,2.1);
  \filldraw[fill=BookTeal!25, draw=BookTeal, line width=0.5pt] (1.715,0) rectangle (2.165,1.7);
  \filldraw[fill=BookTeal!25, draw=BookTeal, line width=0.5pt] (2.335,0) rectangle (2.785,1.3);
  \filldraw[fill=BookTeal!25, draw=BookTeal, line width=0.5pt] (2.955,0) rectangle (3.405,0.9);
  \filldraw[fill=BookTeal!25, draw=BookTeal, line width=0.5pt] (3.575,0) rectangle (4.025,0.45);
  \filldraw[fill=BookAmber!20, draw=BookAmber, line width=0.5pt] (4.195,0) rectangle (4.645,0.12);
  \draw[BookRed, line width=0.5pt, dashed] (0.4,2.25) -- (4.7,2.25);
  \node[hw label, anchor=west, text=BookRed] at (4.78,2.25) {P/E 预算};
  \node[hw label, anchor=north] at (0.7,-0.12) {热};
  \node[hw label, anchor=north] at (4.42,-0.12) {冷};
  \node[hw label] at (2.55,-0.75) {新写入总挑年轻块：热块反复被分配};
  \node[hw label] at (2.55,-1.35) {冷块擦除次数恒为零，少数块先撞预算};
  % 右：加静态均衡
  \node[hw label] at (8.35,2.85) {加静态均衡：擦除次数摊平};
  \draw[BookMuted, line width=0.6pt] (6.2,0) -- (10.5,0);
  \filldraw[fill=BookTeal!25, draw=BookTeal, line width=0.5pt] (6.275,0) rectangle (6.725,1.55);
  \filldraw[fill=BookTeal!25, draw=BookTeal, line width=0.5pt] (6.895,0) rectangle (7.345,1.5);
  \filldraw[fill=BookTeal!25, draw=BookTeal, line width=0.5pt] (7.515,0) rectangle (7.965,1.5);
  \filldraw[fill=BookTeal!25, draw=BookTeal, line width=0.5pt] (8.135,0) rectangle (8.585,1.5);
  \filldraw[fill=BookTeal!25, draw=BookTeal, line width=0.5pt] (8.755,0) rectangle (9.205,1.5);
  \filldraw[fill=BookTeal!25, draw=BookTeal, line width=0.5pt] (9.375,0) rectangle (9.825,1.45);
  \filldraw[fill=BookTeal!25, draw=BookTeal, line width=0.5pt] (9.995,0) rectangle (10.445,1.45);
  % 冷数据搬向高计数块
  \draw[hw return] (10.22,1.7) -- (10.22,2.45) -- (6.6,2.45) -- (6.6,1.7);
  \node[hw label, anchor=south] at (8.4,2.5) {冷数据搬进高计数块};
  \node[hw label] at (8.35,-0.75) {年轻块换出来承接新写，各块摊平};
  \node[hw label] at (8.35,-1.35) {均衡搬运本身计入写放大};
\end{pdfdiagram}
\diagramnote{磨损均衡前后的擦除次数分布（块按擦除次数排列）。只有动态均衡时，新写入总挑擦除次数最少的块，热数据块反复被分配、先撞到 P/E 预算，而冷数据占着的块擦除次数恒为零——少数块先擦穿，整盘提前报废。静态磨损均衡定期把长期不改写的冷数据搬进擦除次数多的块里"养老"，把年轻块换出来轮换，各块擦除次数被摊平；均衡搬运本身计入写放大，寿命的公平是花钱买的。}

盘厂把寿命预期做成一个标称值：TBW（total bytes
written），承诺"累计主机写入达到此数之前，盘在规格内工
作"。它按主机写入计、不按闪存实际写入计——写放大是盘
厂自己的成本，TBW 是给用户的承诺。拿 TBW 反推每日预算
是一道除法。例：256 GB 消费盘，TBW 150 TB，每天写
40 GB：

\begin{lstlisting}
寿命 = 150 TB / (40 GB/天) = 150000 / 40 = 3750 天，约 10.3 年
\end{lstlisting}

约十年，超过多数机器的服役期——前提写在算式背后：WA
低、负载里没有集中的擦写热点。同一张标称，不同的每日写
入量对应完全不同的寿命预期：

\begin{longtable}{L{0.2\linewidth}L{0.2\linewidth}L{0.2\linewidth}}
\toprule
每天写入量 & 预算天数 & 折合年限 \\
\midrule
20 GB & 7500 天 & 约 20.5 年 \\
40 GB & 3750 天 & 约 10.3 年 \\
100 GB & 1500 天 & 约 4.1 年 \\
500 GB & 300 天 & 约 0.8 年 \\
\bottomrule
\end{longtable}

同一道题换个角度更有意思：150 TB 相当于把 256 GB 的盘写
满约 $150\,000/256\approx586$ 次，而 TLC 的 P/E 预算还有
1000-3000 次——标称值留出的安全边，正是给写放大与均衡
搬运预备的。若负载让 WA 长期停在 3，同样的主机写入量下，
闪存实际承受约 $586\times3\approx1760$ 次全盘擦写，落在
TLC 预算中段。这也让"TBW 相同的两块盘，体质可能差几倍"
有了定量含义：标称一样，WA 不同，介质实际寿命就不同。

规格书里的配套指标 DWPD（drive writes per day，保修期内
每天可把整盘写几遍）与 TBW 是同一笔账的两种记法：上例每
天 40 GB 合 $40/256\approx0.16$ DWPD，五年保修累计约
$40/256\times365\times5\approx285$ 次全盘写入，仍在 586 次
标称之内。读寿命规格时先分清它在报哪一个，再换算成自己
的每日写入量，比直接比数字可靠。温度是另一个隐形变量：
电荷保持力随温度升高而下降，消费盘的保持力标称大致是
"断电后 30 摄氏度存放一年"；FTL 的定期巡检刷新正是为这
条曲线兜底，刷新本身也耗 P/E，又一笔主机看不见的磨损。

\topic{TRIM 与稳态性能}
最后一个机制同时解释"盘为什么越用越慢"和"为什么能慢得
不那么狠"。文件系统删除文件时只改自己的元数据，从不通知
盘——机械盘时代这是美德，老数据原地留着，误删还能恢复。
对 SSD 这却意味着两个世界脱节：文件早删了，FTL 还当那些
页是有效数据，GC 照样把它们当宝贝搬走，写放大白白上涨。
盘越用越满，$v$ 越高，WA 沿 $1/(1-v)$ 爬升，随机写越来
越慢——"盘满变慢"的全部机理就在此。

TRIM 是补上脱节的命令：文件系统删文件、释放区间时，顺手
给盘发一条 TRIM（ATA 里是 DATA SET MANAGEMENT 命令的
TRIM 属性，NVMe 里是 Dataset Management 的 deallocate），
告知这段 LBA 不再使用。FTL 收到后直接把这些页标记无效，
下次 GC 它们与从没写过一样，不必再搬。TRIM 不改变任何物
理规律，只改变 FTL 的信息量：$v$ 从此按真实值算，而不是
按"文件系统有史以来写过的所有页"算。顺带一提，被 TRIM
的页读回什么由盘决定（不少盘返回全零），指望删除后还能
恢复文件的软件从此失业。下发时机有两派：挂载选项
\code{discard} 每次删除即时下发，简单但每次删除多一笔命
令开销；\code{fstrim} 定时任务攒一批区间一起发，是多数
发行版的默认。

把出厂态与稳态摆在一起，差距一眼见底：

\begin{longtable}{L{0.26\linewidth}L{0.2\linewidth}L{0.14\linewidth}L{0.26\linewidth}}
\toprule
状态 & 空块来源 & WA 量级 & 4 KiB 随机写 IOPS 量级 \\
\midrule
出厂态（新盘） & 全部块空闲 & 1 & 数万 \\
稳态，有 TRIM 与预留 & GC 回收 & 2-3 & 万级 \\
稳态，无 TRIM 且盘满 & GC 反复搬运 & 5-10 & 千级，伴随时延尖峰 \\
\bottomrule
\end{longtable}

新盘所有块皆空，写请求直接编程空页，WA = 1；持续随机写
把盘推入稳态后，每次写都可能触发 GC，IOPS 按 WA 的倒数
缩水，时延还伴随 GC 尖峰——主机在队列那头看到的现象、
它对本章第五节尾延迟的含义，留到那时再算。企业盘规格书
分报"出厂态"与"稳态"两组随机写数字，正因为它们本就是两
个量；消费盘多半只报前者，读规格书时这是第一个要问的问
题。要把整块盘拉回出厂态，手段是全盘擦除：ATA 的
\code{SECURITY ERASE}、NVMe 的 Format/Sanitize，或对所有
已用区间一次性 TRIM——映射表清空、块全部回收，WA 回到
1；二手交易前的数据清除与性能恢复，是同一条命令的两份
用途。

稳态性能还有队列一侧的条件，用第十章的 Little 定律随手
可算：4 KiB 随机读时延约 80 $\mu$s，队列深度 1 时吞吐上
限是 $1/80\ \mu\mathrm{s}=12\,500$ IOPS，约 50 MB/s；想
吃到 NVMe 标称的 50 万 IOPS，需要随时有
$5\times10^{5}\times80\ \mu\mathrm{s}=40$ 个请求在途。这
就是深队列存在的理由，也是"单线程跑不满 SSD"的正解——
不是盘虚标，是队列深度不够。维护一块盘的长期性能，说到
底三件事：留足空闲容量压低 $v$、打开 TRIM 让 FTL 知情、
别把随机写当顺序写使。三条都不花钱，花的只是对介质脾气
的理解。

\section{三、机械硬盘与访问调度}

机械盘是主存之外仍在大量服役的机电设备：存储介质是磁，定位机构
是音圈电机加主轴电机。它的全部性能特征可以归结为一个事实——盘
片上只有两个运动部件：摆动的磁臂和旋转的盘片，两者的时间常数都
是毫秒级；而读通道、盘上缓存、接口电路这些电的部分都在纳秒到微
秒级。机械盘几十年的性能工程，本质上是一部"怎样躲开这两件机械
运动"的历史。

\topic{磁道、柱面与扇区}盘片是铝合金或玻璃基板，两面涂磁性材料，
每面配一个磁头；所有磁头固定在同一根磁臂上同步摆动，盘片由主轴
电机带着恒速旋转。磁头停在某个半径上不动、盘片转过一周，磁头在
盘面上扫出的圆环叫磁道；一张盘面上有几十万条同心磁道。所有盘面
上同一半径的磁道集合叫柱面——这些圆环在空间里恰好叠成一个圆柱
面。寻道完成后整组磁头同时就位，所以连续数据按柱面组织：在同一
柱面内换面读不用再挪磁臂，只需电子开关换一个磁头。磁道沿圆周切
成小段，每段是一个扇区：传统容量 512 字节，近年的 Advanced
Format 把物理扇区改为 4096 字节，同时以 512e 模式向旧软件模拟
512 字节扇区。

先建立结构上的数量级直觉。一块 7200rpm 的 4TB 桌面盘通常用三张
碟片、六个磁头，单碟约 1.3TB；每张盘面刻下几十万条磁道，道间距
以十纳米计；磁头的飞行高度只有几纳米，比烟尘颗粒还小几个数量级
——这就是盘体必须密封、怕震、一粒灰就可能划伤盘面的原因。机械
盘的全部精密性都花在"让磁头以几纳米高度贴着盘面飞、又永远不许
碰上去"这一件事上。

磁记录方式本身也换过几代。纵向磁记录（LMR）把磁畴平铺在盘面内，
面密度很快撞上极限；垂直磁记录（PMR）把磁畴竖起来写，配合更窄的
读磁头，面密度提高了一个量级，今天的 CMR 盘都是 PMR。再往后，写
磁头的宽度做不下去了——磁场不能无限聚焦——SMR 干脆利用"写宽
读窄"的差值让磁道重叠，这是本节后段的主题；HAMR 与 MAMR 则用激
光或微波辅助写入更小的晶粒，仍在爬坡。磁密度的每一步都同时改写
容量、速度与"介质的脾气"，这正是本章要一类设备一节分开讲的原
因。

\begin{longtable}{L{0.14\linewidth}L{0.40\linewidth}L{0.34\linewidth}}
\toprule
结构要素 & 物理实质 & 典型数量级 \\
\midrule
盘片 platter & 铝合金或玻璃基板，双面涂磁，恒速旋转 & 7200rpm 桌面盘 1--4 张碟 \\
磁头 head & 每面一个，固定在磁臂上由音圈电机驱动 & 飞行高度仅数纳米 \\
磁道 track & 磁头不动、盘转一周扫出的圆环 & 每面几十万条，道间距数十纳米 \\
柱面 cylinder & 所有盘面同半径磁道的集合 & 柱面内换面不再寻道 \\
扇区 sector & 磁道沿圆周的最小读写单位 & 512 字节或 4096 字节 \\
\bottomrule
\end{longtable}

区域位记录（ZBR）决定了一块盘的容量与速度分布。早期磁盘每条磁
道的扇区数相同，容量被最内圈锁死：外圈周长是内圈两倍多，却写同
样多的位，外圈的线密度被白白浪费。区域位记录把盘面按半径分成十
几个区，区内每条磁道的扇区数固定，越靠外的区每道扇区越多。算一
笔几何账：数据区外半径约 43mm、内半径约 20mm，周长比约
$43/20\approx2.15$，恒定线密度下外圈每条磁道能容纳的扇区数约为
内圈的 2 倍。转速恒定，单位时间从磁头下流过的字节数正比于每道
扇区数，于是外圈的顺序读速约为内圈两倍：同一块盘，外圈约
200MB/s，扫到内圈只剩约 100MB/s。LBA 编号从外圈排起，"分区越靠
前越快"这条老经验，物理根源就在 ZBR。

4K 物理扇区带来一条连带约定：分区与文件系统必须按 4K 对齐。若文
件系统的 4K 块错开一个 512 字节逻辑扇区，每一次块写都会横跨两个
物理扇区，盘内要先读出两个 4K、合并、再写回——一次写放大成两
次读加两次写。512e 模式把这件事藏进固件，性能账却藏不住；4Kn 盘
干脆只暴露 4K 扇区，要求软件直接按物理现实对齐。

扇区的传统地址是三元组（柱面， 磁头， 扇区）——这三个数曾经是货
真价实的几何坐标：柱面号指挥磁臂停到哪条半径，磁头号选中哪张盘
面，扇区号等盘片转到第几段。第六章 int 13h 的寄存器约定
\code{CH}、\code{CL}、\code{DH} 就按这三个量命名；BIOS 上限 1024
柱面、255 磁头、63 扇区乘 512 字节，拼出约 8GB 的寻址上限，正是
几何时代的遗产。ZBR 与坏道重映射普及之后，盘内真实布局对外不再
透明，CHS 退化成一组虚构参数，BIOS 与 ATA 各自做一层"假几何"翻
译；LBA 随后把寻址拍平成从 0 开始的扇区序号，盘固件收到 LBA 后
自己换算（区， 道， 扇区）。磁盘地址从几何坐标变成逻辑序号，是外
设把物理复杂性藏进固件的典型一例，也是本书反复出现的主题：约定
一旦立下，硬件在约定后面换了几轮实现，约定本身仍在服役。

\topic{寻道与旋转延迟}一次扇区访问的时间拆成三段：寻道
（seek），磁臂移到目标磁道并安定下来；旋转等待（rotational
latency），等目标扇区转到磁头下方；传输，数据经读通道、盘上缓
存、接口由 DMA 搬进内存——最后这段走的就是第五章描述符与 DMA
的通道，微秒级。前两段是机械运动，毫秒级。随机小请求的总耗时几
乎全在前两段，这是本节的中心事实。

把三段耗时按真实比例画在一条时间条上，与盘片的结构对照
着看：

\begin{pdfdiagram}
  % 左：盘片俯视——磁道、扇区、磁臂
  \draw[BookMuted, line width=0.55pt] (1.8,0.6) circle (1.5);
  \draw[BookMuted, line width=0.45pt] (1.8,0.6) circle (1.05);
  \draw[BookMuted, line width=0.45pt] (1.8,0.6) circle (0.6);
  \fill[BookTeal!20] ([shift={(1.8,0.6)}]205:0.6) arc[start angle=205, end angle=235, radius=0.6] -- ([shift={(1.8,0.6)}]235:1.5) arc[start angle=235, end angle=205, radius=1.5] -- cycle;
  \draw[BookMuted, line width=0.5pt] (0.7,-0.5) -- (0.7,-1.05);
  \node[hw label] at (0.7,-1.28) {扇区（512 B / 4 KB）};
  \draw[BookMuted, line width=0.5pt] (1.8,2.1) -- (1.8,2.4);
  \node[hw label, anchor=south] at (1.8,2.42) {磁道（每面数十万条同心环）};
  \draw[BookMuted, line width=2.0pt] (4.6,-0.1) -- (3.05,-0.1);
  \fill[BookMuted] (4.6,-0.1) circle (0.12);
  \node[hw state, minimum width=0.3cm, minimum height=0.3cm, inner sep=1pt] at (2.95,-0.1) {};
  \draw[BookMuted, line width=0.5pt] (4.6,-0.24) -- (4.6,-0.5);
  \node[hw label, anchor=north] at (4.6,-0.55) {磁臂与磁头（音圈电机驱动）};
  \node[hw label] at (1.8,-1.8) {盘片恒速旋转（如 7200 rpm），磁头等目标扇区转到位};
  % 右：一次随机 4 KB 读的时间条
  \node[hw label, anchor=west] at (5.3,2.85) {随机 4 KB 读一次，合计 $\approx$9.2 ms};
  \filldraw[fill=BookAccent!18, draw=BookAccent, line width=0.5pt] (5.3,1.0) rectangle (7.9,1.7);
  \node[hw label, text=BookInk] at (6.6,1.35) {寻道 $\approx$5 ms};
  \filldraw[fill=BookAmber!20, draw=BookAmber, line width=0.5pt] (7.9,1.0) rectangle (10.3,1.7);
  \node[hw label, text=BookInk] at (9.1,1.35) {旋转等待 $\approx$4.17 ms};
  \filldraw[fill=BookTeal!30, draw=BookTeal, line width=0.5pt] (10.3,1.0) rectangle (10.44,1.7);
  \draw[BookMuted, line width=0.5pt] (10.37,1.7) -- (10.37,2.2);
  \node[hw label, anchor=south] at (10.37,2.22) {传输 $\approx$0.02 ms};
  \draw[BookMuted, line width=0.6pt, <->] (5.3,0.78) -- (10.3,0.78);
  \node[hw label] at (7.8,0.5) {机械运动，毫秒级，占约 99.8\%};
  \node[hw label, anchor=west] at (5.3,0.05) {IOPS $\approx$ 1000/9.2 $\approx$ 109（"约 100 IOPS"的由来）};
\end{pdfdiagram}
\diagramnote{机械盘的结构与一次随机 4 KB 读的时间构成。左图俯视盘片：同心圆环是磁道，圆周上的一小段是扇区，磁臂把磁头送到目标半径后，等目标扇区转到磁头下方。右图把三段耗时按真实比例排开：寻道约 5 ms、旋转等待约 4.17 ms 都是毫秒级的机械运动，传输约 0.02 ms 是微秒级的电过程——随机小请求的时间几乎全部花在前两段，"一块机械盘约 100 IOPS"由此而来。}

寻道本身还可再拆：磁臂加速、巡航、减速、在目标磁道上安定到可读
写的精度。短寻道花不完加速段，长寻道大部分时间在巡航，所以"平
均寻道时间"必须注明工作负载：厂商标称的 4--10ms 通常指随机访问
下的平均；相邻磁道间的一步只需不到 1ms，从盘片最外到最内的全程
寻道则要 15--20ms。后文电梯调度能省时间，靠的正是把随机寻道换
成道间寻道。

\begin{longtable}{L{0.16\linewidth}L{0.30\linewidth}L{0.42\linewidth}}
\toprule
阶段 & 物理动作 & 典型耗时与决定因素 \\
\midrule
寻道 & 磁臂移动并安定 & 随机平均 4--10ms；道间 <1ms；全程 15--20ms \\
旋转等待 & 等目标扇区转到位 & 平均半圈：7200rpm 约 4.17ms，15000rpm 约 2ms \\
传输 & 读通道、缓存、DMA 进内存 & 4KB 约 0.02ms，由介质读速决定 \\
\bottomrule
\end{longtable}

旋转等待的演算：7200rpm 即每秒 120 转，一转 $60000/7200\approx
8.33$ms；目标扇区等概率出现在圆周任意位置，平均等半圈，即约
4.17ms。这个数由转速唯一决定，软件动不了它；15000rpm 的企业盘
把它减半到 2ms，代价是更小的盘片、更小的容量与更高的功耗。调度
能做的只有重排——让每次定位更接近"顺路"，而不是消灭等待本身。

随机 4KB 读的总账：取随机平均寻道 5ms（4--10ms 区间的中值），加
旋转等待 4.17ms，加传输 0.02ms，合计约 9.2ms，常说"约 10ms"。于
是单盘纯随机 4KB 读每秒最多完成约 $1000/9.2\approx109$ 次，即约
100 IOPS；乘 4KB，吞吐约 0.45MB/s。"一块机械盘约 100 IOPS"这条
经验值就是这么来的：IOPS 不是接口参数，而是机械时间常数的倒数。
工程含义很直接：凡把机械盘当随机小请求设备用的设计——数据库直
接随机写裸盘、把虚机镜像打散在机械盘上——容量规划都要从每盘约
100 IOPS 起步，而不是从标称的 200MB/s 起步。

把这条换算画成计算链，从转速到吞吐的每一步都能复算：

\begin{pdfdiagram}
  % 三段耗时
  \node[hw compute, minimum width=2.3cm] (seek) at (1.65,0) {随机寻道 $\approx$5 ms\\（4--10 ms 中值）};
  \node[hw compute, minimum width=2.3cm] (rot) at (4.35,0) {旋转等待 $\approx$4.17 ms\\（7200 rpm 平均半圈）};
  \node[hw compute, minimum width=2.3cm] (xfer) at (7.05,0) {传输 $\approx$0.02 ms\\（4 KB 经 DMA）};
  \node[hw control, minimum width=6.9cm] (sum) at (4.35,-1.5) {一次随机 4 KB 读 $\approx$ 9.2 ms（机械运动占约 99.8\%）};
  \draw[hw bus] (seek.south) -- (seek.south |- sum.north);
  \draw[hw bus] (rot.south) -- (rot.south |- sum.north);
  \draw[hw bus] (xfer.south) -- (xfer.south |- sum.north);
  \node[hw state, minimum width=5.2cm] (iops) at (3.15,-2.9) {IOPS $=1000/9.2\approx109$（常说约 100）};
  \node[hw memory, minimum width=4.2cm] (bw) at (8.15,-2.9) {吞吐 $=109\times4$ KB $\approx0.45$ MB/s};
  \draw[hw bus] (sum.south) -- (sum.south |- iops.north);
  \draw[hw bus] (iops.east) -- (bw.west);
  \node[hw label] at (4.35,-3.95) {IOPS 不是接口参数，而是机械时间常数的倒数；对照顺序读外圈约 200 MB/s};
\end{pdfdiagram}
\diagramnote{随机 4 KB 读的 IOPS 换算链：随机寻道约 5 ms 加平均半圈的旋转等待约 4.17 ms，再加 4 KB 传输约 0.02 ms，合计约 9.2 ms；IOPS 是它的倒数，约 109，乘 4 KB 得吞吐约 0.45 MB/s——与顺序读外圈约 200 MB/s 差两个数量级。容量规划要从每盘约 100 IOPS 起步，而不是从标称带宽起步。}

短行程（short-stroking）是用容量换定位时间的经典手段：只使用外
圈一小部分柱面，寻道范围被限制在最窄的一弧，平均寻道从 5ms 压到
1--2ms，旋转等待不变，单次随机访问降到约 6ms，IOPS 提到约 170；
外圈同时又是 ZBR 里最快的一区，顺序吞吐也最高。代价是只用盘容
量的一小角。它说明一个事实：机械盘的容量、时延、IOPS 三者可以
互相兑换，兑换率由机械时间常数决定。

这个换算直接进容量规划。一个需要 5000 随机 IOPS 的应用，用机械盘
做 RAID0 需要约 50 块盘（每盘约 100 IOPS），哪怕总容量只需要其
中十块盘的量；同一负载换 SATA SSD（随机读数万 IOPS），一块就
够。存储选型先定 IOPS 指标、再定容量指标，顺序弄反是机械盘时代
最常见的预算事故之一——盘买够了容量，性能却只有预算值的几十分
之一。第十章的 Amdahl 定律在这里换成一句大白话：整个阵列的随机
吞吐，等于盘数乘以每盘那约 100 次。

\topic{顺序与随机的吞吐差距}顺序读时磁头定位一次，之后盘片每转
一圈交出一条完整磁道，接口与缓存全速运行：外圈约 200MB/s，整盘
平均约 150--160MB/s。顺序吞吐由位密度与转速决定，是"电"的速
度；随机吞吐由寻道与旋转决定，是"机械"的速度。同一块盘有两个
速度，差多少完全取决于访问模式。接口在这里不是瓶颈：SATA 6Gb/s
按 8b/10b 编码折算有效带宽约 600MB/s（第八章讨论过这类编码开销
的来源），介质供数最多 200MB/s，线路一直空着一多半；是 SSD 的出
现才把 SATA 真正打满。

三笔演算把差距钉死。吞吐比：$200\div0.45\approx450$ 倍，超过两
个数量级。任务时间：顺序读 1GB 约 5 秒；随机读同样 1GB 是
$2^{30}/2^{12}=262144$ 次 4KB 读，乘 9.2ms 约 2410 秒——5 秒对
40 分钟，读的是同一份数据。整盘扫描：4TB 盘以平均 160MB/s 顺序
扫一遍约 $4\times10^{12}/(1.6\times10^{8})=25000$ 秒，约 7 小时
——阵列重建与备份窗口的时长下限就是这么估的。

这三个数不必靠相信，一台机器几条命令就能复现。fio 的顺序读、队
列深 1 的随机读、队列深 32 的随机读，恰好对应本节的三个吞吐档
位；iodepth 就是在途请求数，与第十章 Little 定律里的在途数是同一
个量。

\begin{lstlisting}
# same 7200rpm disk, three access patterns (Linux, fio)
fio --name=seq   --rw=read     --bs=1M --size=4G --filename=/dev/sdX
#     => about 200 MB/s on outer tracks
fio --name=rand  --rw=randread --bs=4k --size=4G --iodepth=1  ...
#     => about 100 IOPS, i.e. about 0.4 MB/s
fio --name=randq --rw=randread --bs=4k --size=4G --iodepth=32 ...
#     => NCQ reorders in the drive: a few hundred IOPS
\end{lstlisting}

\begin{longtable}{L{0.22\linewidth}L{0.30\linewidth}L{0.36\linewidth}}
\toprule
访问模式 & 每次定位开销 & 7200rpm 盘的吞吐 \\
\midrule
顺序读 & 一次寻道，之后免定位 & 外圈约 200MB/s，平均约 160MB/s \\
随机 4KB，队列深 1 & 每次都寻道加旋转 & 约 100 IOPS，约 0.45MB/s \\
随机 4KB，排队重排后 & 短寻道加部分旋转 & 数百 IOPS，约 1--1.6MB/s \\
\bottomrule
\end{longtable}

物理根源是惯性。磁臂和盘片有质量，加速、减速、安定都要毫秒；而
磁性状态的读出快得多——磁头一旦就位，读通道以每纳秒若干位的速
度采样流过的弧段。"机械盘怕随机"不是介质慢，是定位慢。由此得
到全部工程推论的共同方向：把随机访问变成顺序访问。日志结构文件
系统先写日志、数据库顺序追加 redo、文件系统按分配组聚集数据，
以及下一段的调度与合并，都是这同一件事在不同层的化身。

\topic{电梯调度与合并}块层队列给了软件选择权：队列里同时压着
32 个请求时，先服务谁是自由的。SCAN（电梯算法）把请求按目标柱
面排序，磁臂从一端扫到另一端、沿途服务，到头反向，像电梯不在无
人楼层停靠；LOOK 不到物理端点就提前折返，C-LOOK 只单向服务、到
头快速回扫。排序换掉的是"随机平均寻道 5ms"这个前提：排序后相
邻请求大多位于相邻柱面，寻道从毫秒级的全程动作退化成不到 1ms
的道间一步。

与排序同时发生的是合并：LBA 相邻的请求在块层拼成一个更大的请
求。8 个相邻 4KB 读合成一次 32KB 读，一次定位全部到手；排序吃的
是磁道维度上的重排，合并吃的是"LBA 相邻本就大概率在同一条磁
道"。请求在队列里停留的时间因此有了双重价值——既等来了排序的
依据，也等来了可合并的邻居。

收益演算：32 个随机 4KB 请求按到来顺序（FCFS）服务，约
$32\times9.2\approx294$ms。电梯排序后，磁臂一趟扫过请求分布的柱
面区间，每请求摊到短寻道加一次旋转等待，按
$0.8+4.17+0.02\approx5$ms 粗算，约 160ms；LBA 相邻的请求再被合
并，总时间更短。注意这个收益的前提：队列里得有东西可排。队列深
度为 1 时没有重排空间，调度无能为力。这正是第十章 Little 定律的
另一面——在途请求数既是吞吐的原料，也是调度的原料；评测随机性
能必须把队列深度写进条件，QD1 的约 100 IOPS 与 QD32 的数百 IOPS
是同一块盘的两种真实，谁也不是"作弊"。

把同一批请求的两种服务顺序画成磁臂轨迹，排序省下的是什么就有了形状：

\begin{pdfdiagram}
  % 横轴：时间（服务顺序）；纵轴：柱面 0-100
  \draw[BookMuted, line width=0.6pt, ->] (0,0) -- (10.4,0);
  \draw[BookMuted, line width=0.6pt, ->] (0,0) -- (0,4.5);
  \node[hw label, anchor=north] at (5.2,-0.6) {时间（服务顺序）};
  \node[hw label, rotate=90, anchor=south] at (-0.85,2.2) {柱面};
  \foreach \y/\t in {0/0, 2/50, 4/100} {
    \draw[BookMuted, line width=0.5pt] (-0.07,\y) -- (0.07,\y);
    \node[hw label, anchor=east] at (-0.12,\y) {\t};
  }
  % FCFS：按到达顺序，来回甩动
  \draw[BookAccent, line width=0.9pt] (0.7,0.8) -- (2.1,3.4) -- (3.5,1.6) -- (4.9,2.4) -- (6.3,0.4) -- (7.7,3.6) -- (9.1,2.2);
  % SCAN：按柱面排序，单趟扫过
  \draw[BookTeal, line width=0.9pt] (0.7,0.4) -- (2.1,0.8) -- (3.5,1.5) -- (4.9,2.2) -- (6.3,2.4) -- (7.7,3.4) -- (9.1,3.6);
  \foreach \p in {(0.7,0.8), (2.1,3.4), (4.9,2.4), (6.3,0.4), (7.7,3.6), (9.1,2.2)} {
    \fill[BookAccent] \p circle (0.05);
  }
  \foreach \p in {(0.7,0.4), (2.1,0.8), (4.9,2.2), (6.3,2.4), (7.7,3.4), (9.1,3.6)} {
    \fill[BookTeal] \p circle (0.05);
  }
  \node[hw label, anchor=west] at (0.2,4.2) {同一批 32 个随机请求};
  \node[hw label, anchor=west] at (0.2,3.75) {省下的寻道时间：约 130 ms};
  \draw[BookAccent, line width=0.9pt] (5.9,4.2) -- (6.9,4.2);
  \node[hw label, anchor=west] at (7.0,4.2) {FCFS：约 294 ms};
  \draw[BookTeal, line width=0.9pt] (5.9,3.75) -- (6.9,3.75);
  \node[hw label, anchor=west] at (7.0,3.75) {SCAN：约 160 ms};
\end{pdfdiagram}
\diagramnote{同一批随机请求的两种磁臂轨迹（示意其中七个请求的位置）。FCFS 按到达先后服务，磁臂在柱面间来回甩动，32 个请求全程约 $32\times9.2\approx294$ ms；SCAN 按柱面排序后单趟扫过，每请求只摊到短寻道加一次旋转等待，全程约 160 ms。两条轨迹之间省下的约 130 ms，全是寻道时间。}

\begin{longtable}{L{0.20\linewidth}L{0.34\linewidth}L{0.34\linewidth}}
\toprule
调度策略 & 服务顺序 & 关键问题 \\
\midrule
FCFS & 按到达先后 & 无重排，随机负载下磁臂来回甩 \\
SATF & 最短寻道优先 & 吞吐最好，远端请求可能饿死 \\
SCAN / LOOK & 按柱面排序，到头折返 & 兼顾公平，折返点附近等待偏长 \\
C-LOOK & 单向扫描，到头回扫 & 等待更均匀，回扫本身空跑一趟 \\
mq-deadline & 排序加截止时间，读优先于写 & 读期限远短于写期限，防止饿死 \\
none & 不调度 & SSD 无机械定位，排序是净开销 \\
\bottomrule
\end{longtable}

读优先于写有应用层的原因：读大多是同步等待，应用线程停在原地；
写走页缓存 write-back，应用写完就走。把写无限推迟却不行——脏
页终究要回写，推迟过久会把抖动挪到更糟的时刻，所以调度器给写也
设期限，只是宽得多。调度器也不是越聪明越好：机械定位消失之后，
排序本身成了净开销，介质换成 SSD，内核调度器就该退场。软件结构
跟着介质的物理常数走，这是外设栈的普遍规律。

期限的具体值可以读出来：Linux 的 mq-deadline 给读请求默认约
500ms、写请求约 5s 的期限，到期未服务的请求优先发出，饿死被时
间上限封死。调度器与队列预算也可以直接观察：

\begin{lstlisting}
$ cat /sys/block/sda/queue/scheduler
none [mq-deadline] kyber bfq      # current: mq-deadline
$ cat /sys/block/sda/queue/nr_requests
128                               # block layer queue depth budget
\end{lstlisting}

块层队列预算（nr\_requests）决定最多能攒多少请求供排序与合并：
太小，重排空间不足；太大，单请求的排队时延上升。同一个旋钮，在
机械盘上和 NVMe 上的合适应答完全不同——第五节会再回到这个数。

测量的另一半在 iostat 里：r\_await 与 w\_await 把每次读写的平均
总时延以毫秒报出，机械盘随机负载下看到 5--10ms，正是寻道加旋转
的直接读数；await 突然涨到几十毫秒，多半意味着队列在某一层堵住
或盘内在做重映射。第十章"用测量代替猜测"的原则在外设上格外省
钱：磁盘几乎从不说谎，它的毫秒数摆在那里，与本节的时间构成一项
一项核对即可。

\topic{NCQ 与 SMR}操作系统的排序发生在驱动里，驱动不知道盘此刻
磁头停在哪、盘片转到什么角度；盘固件知道。SATA 的原生命令队列
（NCQ）允许主机一次提交最多 32 条命令，盘内按"当前柱面加旋转相
位"重排：同一柱面的多个请求按扇区转到磁头下的先后执行，旋转等
待从平均半圈降到更小。同样是 32 个随机 4KB 请求，盘内重排后每条
摊到的定位时间约 2.5--4ms，IOPS 从约 100 提到约 250--400。这正
解释了"一次随机读约 10ms"与"标称数百 IOPS"为何同时成立：前者
是单请求的物理时延，后者是排队重排后的吞吐；时延没变，在途数变
了，按第十章 Little 定律，吞吐随之改变。AHCI 只有一条 32 深的队
列；更深的队列模型——NVMe 的 64K 条队列、每条 64K 深——为 SSD
准备，本章第五节展开。对机械盘，32 个在途请求提供的重排空间已
把大部分收益吃掉。

重排做的事用六个请求示意便清楚（实际队列里是 32 个）：

\begin{pdfdiagram}
  % 上：到达顺序，相位乱跳
  \node[hw label, anchor=east] at (0.55,0) {到达顺序};
  \node[hw compute, minimum width=1.1cm, minimum height=0.6cm, inner sep=1pt] (a1) at (1.6,0) {相位 5};
  \node[hw compute, minimum width=1.1cm, minimum height=0.6cm, inner sep=1pt] (a2) at (3.15,0) {相位 1};
  \node[hw compute, minimum width=1.1cm, minimum height=0.6cm, inner sep=1pt] (a3) at (4.7,0) {相位 6};
  \node[hw compute, minimum width=1.1cm, minimum height=0.6cm, inner sep=1pt] (a4) at (6.25,0) {相位 2};
  \node[hw compute, minimum width=1.1cm, minimum height=0.6cm, inner sep=1pt] (a5) at (7.8,0) {相位 4};
  \node[hw compute, minimum width=1.1cm, minimum height=0.6cm, inner sep=1pt] (a6) at (9.35,0) {相位 3};
  \draw[hw bus] (a1.east) -- (a2.west);
  \draw[hw bus] (a2.east) -- (a3.west);
  \draw[hw bus] (a3.east) -- (a4.west);
  \draw[hw bus] (a4.east) -- (a5.west);
  \draw[hw bus] (a5.east) -- (a6.west);
  \node[hw label] at (5.5,-0.82) {按到达先后服务：寻道加平均半圈等待，每条约 10 ms $\Rightarrow$ 约 100 IOPS};
  % 下：按旋转相位重排
  \node[hw label, anchor=east] at (0.55,-1.9) {NCQ 重排后};
  \node[hw compute, minimum width=1.1cm, minimum height=0.6cm, inner sep=1pt] (b1) at (1.6,-1.9) {相位 1};
  \node[hw compute, minimum width=1.1cm, minimum height=0.6cm, inner sep=1pt] (b2) at (3.15,-1.9) {相位 2};
  \node[hw compute, minimum width=1.1cm, minimum height=0.6cm, inner sep=1pt] (b3) at (4.7,-1.9) {相位 3};
  \node[hw compute, minimum width=1.1cm, minimum height=0.6cm, inner sep=1pt] (b4) at (6.25,-1.9) {相位 4};
  \node[hw compute, minimum width=1.1cm, minimum height=0.6cm, inner sep=1pt] (b5) at (7.8,-1.9) {相位 5};
  \node[hw compute, minimum width=1.1cm, minimum height=0.6cm, inner sep=1pt] (b6) at (9.35,-1.9) {相位 6};
  \draw[hw bus] (b1.east) -- (b2.west);
  \draw[hw bus] (b2.east) -- (b3.west);
  \draw[hw bus] (b3.east) -- (b4.west);
  \draw[hw bus] (b4.east) -- (b5.west);
  \draw[hw bus] (b5.east) -- (b6.west);
  \node[hw label] at (5.5,-2.72) {按转到磁头下的先后执行：每条摊定位约 2.5--4 ms $\Rightarrow$ 约 250--400 IOPS};
\end{pdfdiagram}
\diagramnote{同一批随机请求的两种服务顺序（示意六个，实际是 32 个在途）。按到达先后服务，相邻请求相位乱跳，每条都要付一次寻道加平均半圈的旋转等待；盘内固件知道磁头位置与盘片角度，按"当前柱面加旋转相位"重排后，同一柱面的请求按转到磁头下的先后执行，旋转等待被摊薄，每条定位约 2.5--4 ms。时延没变，在途数变了——IOPS 从约 100 提到约 250--400。}

SMR（叠瓦式磁记录）从另一头榨容量。写入磁道天生比读取所需宽度
宽，叠瓦布局让后一条磁道盖住前一条的边缘，像屋顶瓦片一样部分重
叠，磁道密度因此提高，单盘容量便宜约两成。代价全在写的方向：写
第 N 道会顺带覆盖第 N+1 道的重叠区，于是一组磁道构成一条带
（band），典型 256MB，带内只能从带首顺序追加；要改写带中间一个
扇区，必须把它之后的内容读出、合并、整条带重写。极端情况下一次
4KB 随机写牵连 256MB 的重写，写放大约六万五千倍；盘内固件当然不
会每次真这么做，它用下一段的缓存机制把账延后，但账不会消失。

叠瓦的结构与改写的牵连画在一起看，"便宜两成"的代价
就具体了：

\begin{pdfdiagram}
  % 五条叠瓦磁道：后一条盖住前一条的边缘
  \draw[BookTeal, line width=0.42pt, fill=white] (0.7,0) rectangle (6.7,0.55);
  \draw[BookTeal, line width=0.42pt, fill=white] (0.7,-0.35) rectangle (6.7,0.2);
  \draw[BookTeal, line width=0.42pt, fill=white] (0.7,-0.7) rectangle (6.7,-0.15);
  \draw[BookTeal, line width=0.42pt, fill=white] (0.7,-1.05) rectangle (6.7,-0.5);
  \draw[BookTeal, line width=0.42pt, fill=white] (0.7,-1.4) rectangle (6.7,-0.85);
  % 重叠区（被后一道盖住的边缘）
  \fill[BookAmber!25] (0.7,0) rectangle (6.7,0.2);
  \fill[BookAmber!25] (0.7,-0.35) rectangle (6.7,-0.15);
  \fill[BookAmber!25] (0.7,-0.7) rectangle (6.7,-0.5);
  \fill[BookAmber!25] (0.7,-1.05) rectangle (6.7,-0.85);
  \node[hw label] at (1.55,0.375) {道 N};
  \node[hw label] at (1.55,0.025) {道 N+1};
  \node[hw label] at (1.55,-0.325) {道 N+2};
  \node[hw label] at (1.55,-0.675) {道 N+3};
  \node[hw label] at (1.55,-1.025) {道 N+4};
  % 改写点：写第 N+2 道会顺带覆盖第 N+3 道的重叠区
  \draw[BookRed, line width=0.5pt, fill=BookRed!30] (3.6,-0.5) rectangle (4.4,-0.15);
  \node[hw label, anchor=west, text=BookRed] at (4.5,-0.325) {改写点};
  % 读/写宽度对照
  \draw[BookMuted, line width=0.6pt, <->] (7.15,0.2) -- (7.15,0.55);
  \node[hw label, anchor=west] at (7.3,0.375) {读所需宽度};
  \draw[BookMuted, line width=0.6pt, <->] (8.85,-0.35) -- (8.85,0.2);
  \node[hw label, anchor=west] at (9.0,-0.075) {写入磁道宽度};
  % 带
  \draw[BookAmber, line width=0.6pt] (0.5,-1.4) -- (0.5,0.55);
  \node[hw label] at (3.7,-2.1) {一条带（band，典型 256 MB）：带内只能从带首顺序追加};
  \node[hw label, text=BookRed] at (3.7,-2.75) {改写带中一个扇区：其后各道读出、合并，整条带重写——极端时 4 KB 写牵连 256 MB};
\end{pdfdiagram}
\diagramnote{SMR 叠瓦磁道的结构（深色条带是被后一道盖住的重叠区）。写入磁道天生比读取所需的宽度宽，叠瓦布局让后一道盖住前一道的边缘，磁道密度提高，单盘容量便宜约两成。代价在写的方向：写第 N 道会顺带覆盖第 N+1 道的重叠区，于是一组磁道构成一条带，带内只能从带首顺序追加；要改写带中间一个扇区，必须把它之后的内容读出、合并、整条带重写——极端情况一次 4 KB 随机写牵连 256 MB 重写，写放大约六万五千倍。}

盘管式 SMR（DM-SMR）用一小片 CMR 缓存区先收下随机写，空闲时再
整理回叠瓦区：轻负载看不出差别，持续随机写打穿缓存之后，稳态写
吞吐可跌到个位数 MB/s，单次写延迟出现秒级尖峰；NAS 阵列重建时
SMR 盘超时掉盘的案例，根源即此。主机管理型（HM-SMR）索性要求主
机按带分段、带内顺序写，把约束摆上接口，换取确定性——这与
NVMe 分区命名空间（ZNS）是同一思想：介质把脾气讲清楚，软件按脾
气写，双方都轻松。

\begin{longtable}{L{0.22\linewidth}L{0.30\linewidth}L{0.36\linewidth}}
\toprule
特征 & CMR 传统盘 & SMR 叠瓦盘 \\
\midrule
磁道布局 & 磁道分开，可原位改写 & 磁道重叠，带内只能顺序追加 \\
随机写稳态 & 约 100--数百 IOPS & 缓存打穿后可跌到个位数 MB/s \\
缓存特征 & 常见 64--256MB & 同容量点常配异常大的缓存 \\
典型定位 & 通用、NAS、监控 & 归档、冷数据、顺序写为主 \\
阵列重建 & 按标称 IOPS 估算时间 & 持续负载下可能超时掉盘 \\
确认途径 & 厂商型号规格页 & 查厂商公布的 SMR 型号清单 \\
\bottomrule
\end{longtable}

SMR 值得单独成题，还因为一段行业插曲：2020 年前后，厂商曾把
DM-SMR 盘不加说明地混进 NAS 与桌面产品线，用户在做阵列重建与持
续随机写时遭遇数小时卡顿乃至掉盘，追查之下才发现同一型号存在
CMR 与 SMR 两个版本。事件以三家大厂公开 SMR 型号清单收场。教训
写在规格页的边角：同一名称的盘，介质组织方式可以不同，而介质组
织方式决定的是写的语义，不只是速度。读规格书时，CMR 还是 SMR
这一栏与容量、转速同等重要——它回答的是"这块盘允许什么样的访
问模式"。

选购识别并不靠猜：三家大厂在 2020 年后都公布了 SMR 型号清单；同
系列里缓存异常大的容量点值得警惕；商品页出现"归档""冷数据"定
位词时按 SMR 理解。用途涉及随机写、数据库、RAID 重建的，明确选
CMR。SMR 不是"坏盘"，它把"只能顺序写"这条介质的脾气换成容量
与价格；买错场景才是问题。本节结尾回看章首那句：让介质按其擅长
的方式工作——对机械盘，"擅长"二字几乎就是"顺序"的同义词。

\section{四、网卡与网络队列}

网卡面对的约束与磁盘相反：磁盘是"你求它快"，网络是"它不管你收
不收"。链路上的帧按线速到达，不能暂停、不可减速；对端何时发、
发多快，本机说了不算。网卡的全部队列设计都在回答同一个问题：帧
已经在线上了，CPU 还没准备好，怎么办。本节沿帧的生命周期走一
遍：线上怎么数帧，帧进内存走哪条环，CPU 以什么节奏被叫醒，多核
怎样分流，以及数据怎样少搬一次是一次。

\topic{帧、MTU 与线速}以太帧在线上的完整样子：先是 7 字节前导码
加 1 字节帧起始符（SFD），让接收方的时钟恢复电路锁定比特边界；
随后 6 字节目的 MAC、6 字节源 MAC、2 字节类型字段（EtherType），
指明载荷交给谁——IPv4、IPv6 还是 ARP；载荷 46--1500 字节，不足
46 字节必须填充；最后 4 字节 FCS，对整帧做 CRC32 校验，错帧在网
卡层面直接丢弃，不进内存。帧与帧之间还有 12 字节的帧间隙
（IFG），是链路上的静默间隔。MTU 1500 指的是载荷上限；线上一帧
实为 14 字节头加 1500 字节载荷加 4 字节 FCS，共 1518 字节，再算
上前导与帧间隙，每帧占线 1538 字节。

把这一帧的字节布局画出来，固定开销与载荷的比例就有了形状：

\begin{pdfdiagram}
  % 线上一帧：字段宽度为示意，总长对应 1538 字节
  \node[hw label, anchor=south west] at (0,3.3) {线上一帧共占 1538 字节（字段宽度为示意）};
  \draw[BookTeal, line width=0.8pt, fill=BookCodeBg] (0,2.3) rectangle (1.0,3.2);
  \node[align=center] at (0.5,2.75) {前导\\8 B};
  \draw[BookTeal, line width=0.8pt, fill=BookCodeBg] (1.0,2.3) rectangle (2.1,3.2);
  \node[align=center] at (1.55,2.75) {目的\\6 B};
  \draw[BookTeal, line width=0.8pt, fill=BookCodeBg] (2.1,2.3) rectangle (3.2,3.2);
  \node[align=center] at (2.65,2.75) {源\\6 B};
  \draw[BookTeal, line width=0.8pt, fill=BookCodeBg] (3.2,2.3) rectangle (4.1,3.2);
  \node[align=center] at (3.65,2.75) {类型\\2 B};
  \draw[BookAccent, line width=0.9pt, fill=BookCodeBg] (4.1,2.3) rectangle (9.0,3.2);
  \node at (6.55,2.75) {载荷 46--1500 B（MTU 1500）};
  \draw[BookTeal, line width=0.8pt, fill=BookCodeBg] (9.0,2.3) rectangle (9.9,3.2);
  \node[align=center] at (9.45,2.75) {FCS\\4 B};
  \draw[BookMuted, line width=0.8pt, fill=white] (9.9,2.3) rectangle (11.0,3.2);
  \node[align=center] at (10.45,2.75) {IFG\\12 B};
  % 开销占比对照：小帧 vs 大帧
  \node[hw label, anchor=east] at (-0.15,1.35) {最小帧 64 B};
  \draw[BookRed, line width=0.8pt, fill=white] (0,1.0) rectangle (2.0,1.7);
  \node[hw label] at (1.0,1.35) {开销 45\%};
  \draw[BookAccent, line width=0.8pt, fill=BookCodeBg] (2.0,1.0) rectangle (4.5,1.7);
  \node[hw label] at (3.25,1.35) {载荷};
  \node[hw label, anchor=west] at (4.7,1.35) {线上 84 B，载荷率约一半（548 Mb/s）};
  \node[hw label, anchor=east] at (-0.15,0.35) {最大帧 1518 B};
  \draw[BookRed, line width=0.8pt, fill=white] (0,0.0) rectangle (0.11,0.7);
  \draw[BookAccent, line width=0.8pt, fill=BookCodeBg] (0.11,0.0) rectangle (4.5,0.7);
  \node[hw label] at (2.3,0.35) {载荷};
  \node[hw label, anchor=west] at (4.7,0.35) {线上 1538 B，开销仅 2.5\%（975 Mb/s）};
\end{pdfdiagram}
\diagramnote{以太帧在线上的字节布局：前导与 SFD 8 字节用于时钟锁定，MAC 地址与类型字段 14 字节，载荷最多 1500 字节，FCS 4 字节校验，帧间还有 12 字节静默。每帧 38 字节的固定开销与帧长无关：64 字节小帧把它摊成 45\% 的负担，1518 字节大帧只付 2.5\%——压力测试用小帧考 pps、吞吐测试用大帧考带宽，原因就在这张图上。}

\begin{longtable}{L{0.20\linewidth}L{0.12\linewidth}L{0.56\linewidth}}
\toprule
字段 & 字节数 & 作用 \\
\midrule
前导码 + SFD & 8 & 时钟恢复与帧定界，属于物理层，不进内存 \\
目的 MAC & 6 & 网卡按它做第一层过滤，不符则直接丢弃 \\
源 MAC & 6 & 标识发送方，交换机按它学习转发表 \\
类型 EtherType & 2 & 指明载荷协议：IPv4、IPv6、ARP 等 \\
载荷 payload & 46--1500 & MTU 所指的部分，不足 46 字节要填充 \\
FCS & 4 & CRC32 校验，错帧在网卡即丢弃 \\
帧间隙 IFG & 12 & 帧间静默，链路占用但无内容 \\
\bottomrule
\end{longtable}

线速的换算从这里开始。1GbE 的线上比特率是 $10^{9}$ bit/s——第八
章讲眼图时给过这个速率的直观：1Gbps 下每个比特只有 1ns 的判决窗
口。帧越小，每秒能过的帧越多，因为每帧都要付 38 字节的固定开销
（前导 8、帧间隙 12、头 14、FCS 4）。两个端点情形：最小帧 64 字
节，线上占 $64+8+12=84$ 字节即 672 bit，1GbE 每秒过
$10^{9}/672\approx1.49\times10^{6}$ 帧，约 1.49 Mpps；最大帧 1518
字节，线上占 1538 字节即 12304 bit，每秒约
$10^{9}/12304\approx81274$ 帧。10GbE 把这两个数都乘 10：约 14.88
Mpps 与约 81.3 万 pps。

\begin{longtable}{L{0.16\linewidth}L{0.20\linewidth}L{0.24\linewidth}L{0.24\linewidth}}
\toprule
帧长（字节） & 线上占用（字节） & 1GbE（pps） & 10GbE（pps） \\
\midrule
64（最小帧） & 84 & 约 149 万 & 约 1488 万 \\
512 & 532 & 约 23.5 万 & 约 235 万 \\
1518（最大帧） & 1538 & 81274 & 约 81.3 万 \\
\bottomrule
\end{longtable}

有效载荷效率值得单独算：小帧线速下，每秒载荷只有
$1.49\times10^{6}\times46\times8\approx548$ Mb/s，约线速的一半；
大帧下 $81274\times1500\times8\approx975$ Mb/s，接近满载。所以
"线速"与"pps"是两个独立的指标：线速（bit/s）是物理层能力，
pps 度量每帧固定成本的总量——网卡 DMA 一次、描述符消耗一个、协
议栈走一遍，都与帧长无关。压力测试用 64 字节小帧考 pps，吞吐测
试用 1518 字节大帧考带宽，两个数字各回答一个问题；选型时先量业
务的平均帧长，再决定该为哪个指标付钱。

帧格式也不是只有这一种长相。802.1Q 在源 MAC 与类型字段之间插入
4 字节标签，携带 VLAN 编号，帧长上限随之变成 1522 字节；QinQ、
MPLS 再往外套。每多一层封装，MTU 的有效载荷就被咬一口：隧道出
口看到的 1500 字节帧里，真正的数据可能只剩 1400 余字节——VPN
与隧道环境里"MTU 要调小"这条运维经验的来源，就是各层封装头的
总长度。

最小帧 64 字节也有来历。早期共享总线式以太网靠冲突检测
（CSMA/CD）工作：发送方必须发够长的时间，才能让冲突信号从总线
最远端传回来，否则冲突发生而发送方浑然不觉。按 10Mb/s 时代的线
缆长度与传播时延算，这个窗口定为 512 比特时，即 64 字节——最短
合法帧由此而来。交换式全双工以太网早已没有冲突，最小帧却作为约
定的一部分留了下来，今天我们数 pps 时仍在为它付费：每帧 38 字节
的固定开销，一半来自这段历史。

MTU 本身也可以调。数据中心内部常用 9000 字节载荷的巨型帧
（jumbo frame）：线上每帧约 9038 字节，1GbE 每秒帧数降到约 1.4
万，每帧固定成本被摊到六倍的载荷上，pps 与中断、协议栈开销同步
下降约六倍。前提是路径上每一跳的 MTU 都得一致放大，否则分片甚
至静默丢包会把收益全部吃回去——巨型帧是全网约定，不是单机参
数。

"网卡不能喊停"严格说有一个例外：以太网流控允许接收方发 PAUSE
帧，请对端暂停发送指定时长。它的本意是防丢包，实际效果是把缓冲
区不足的问题传染给上游——一台主机收不动，整条链路上的无辜流量
一起被压。现代数据中心普遍关闭全局 PAUSE，丢包问题留给更深的
环、更快的消费和拥塞协议去解决；存储网络使用的按类暂停（PFC）
是另一套权衡，以队头阻塞为代价换无损。这个例外恰好坐实本节的出
发点：默认情况下，到达速率不由接收方控制，缓冲深度与消费速率必
须自己兜底。

\topic{接收环与发送环}网卡与内存之间走的就是第五章的描述符环，
机制原封不动：环是内存里的描述符数组，CPU 推进软件指针、设备推
进硬件指针，doorbell 只提示"指针动了"，完成以设备写回的状态位
为准。网卡的特殊性在于，收发两个方向的环各有一条不对称的约定：
接收环要预投，发送环要回收。

RX 环是预投的。收帧不可预约——对端随时发，网卡不能喊停——所
以驱动必须提前把一批空缓冲区的地址铺满接收环。帧到达后，网卡
MAC 校验 FCS，由 DMA 把帧直接写进下一个空缓冲区，写回描述符状
态（帧长、校验结果、RSS 队列号），再按中断策略通知 CPU。驱动取
走数据、把缓冲区交给协议栈之后，必须立刻补投新的空缓冲区：预投
是义务，不是优化。环上可用缓冲区一旦耗尽，再到的帧没有地方放，
只能丢弃——帧本身没有错，是没人接。

TX 环是回收的。发送由软件发起：驱动把待发帧的描述符挂上发送
环、敲 doorbell；网卡按描述符 gather 数据、送上线路，发完写回完
成状态；驱动在中断或轮询里回收描述符并释放缓冲区。回收不及时，
环被挂满，新帧无处提交，协议栈只能排队等待，背压沿协议栈向上传
导。两个方向的环因此有不同的"满"的含义：RX 环空是丢包，TX 环
满是阻塞。

发送环的回收节奏同样可调：每发完一帧就中断一次太奢侈，驱动通常
让网卡攒一批完成再通知，或干脆在下次发送与定时轮询时顺手回收。
回收延迟直接占用环深——一批未回收的描述符，与一串未补投的空缓
冲区，是同一个问题的两面：环的有效深度，等于标称深度减去滞留在
"已完成未回收"状态里的那部分。

环深与丢包的关系可以直接算。RX 环的深度必须吸收"最坏通知延
迟"内到达的全部帧：1GbE 小帧每秒约 149 万帧，若中断聚合加调度
把最坏通知延迟拖到 200µs，期间到达
$1.49\times10^{6}\times200\times10^{-6}\approx298$ 帧——环深 64
必然丢，512 才够；10GbE 同样的窗口要容纳约 2980 帧，常见取
4096。丢包在以太网层面不报错、不重传，只体现为网卡计数器里的数
字，最后由 TCP 重传表现为"网速偶尔抖一下"。第五章"缓冲区何时
归 CPU、何时归设备"的所有权问题，在这里直接变成丢包与否的分界
线。

\begin{lstlisting}
$ ethtool -g eth0      # show ring depths: RX 512, TX 512
$ ethtool -S eth0 | grep -E 'rx_missed|rx_no_buffer'
     rx_missed_errors: 1382     # frames arrived with no buffer
$ ethtool -G eth0 rx 4096      # deepen RX ring (driver permitting)
\end{lstlisting}

丢包计数是这条链上最诚实的仪表。应用看到的是吞吐波动与重传，协
议栈看到的是乱序，只有网卡计数器直接说出"哪一环没接住"：
rx\_no\_buffer 指预投断供，rx\_missed 指环深或处理速度不够。排障
的顺序因此固定：先看计数器分辨丢在哪一层，再回头调环深、中断参
数或核数——本节后面三个话题，正好各管其中一环。

\topic{中断聚合}每帧一中断在小帧线速下根本不成立。1.49 Mpps 意味
着每 0.67µs 一次中断；一次完整的中断处理——保存现场、读状态、
递交协议栈、恢复现场——按保守的 3µs 计，单核每秒最多服务约 33
万次。需要的算力是 $1.49\times10^{6}\times3\mu s\approx4.5$ 颗核
的全部时间，而且全花在中断进出上，协议栈和应用一行代码都轮不
到：越忙越处理不完，吞吐反而崩，这就是中断风暴（receive
livelock）。0.67µs 这个帧间隔值得记住：它就是"软件为每帧固定成
本付费"的时钟节拍。

把帧间隔与处理时长画在同一条时间轴上，风暴的自我放大就
清楚了：

\begin{waveform}[x=0.8cm,y=0.9cm]
  % 帧到达：每 0.67 µs 一帧
  \foreach \x in {1,2,...,12} \draw[BookTeal, line width=0.9pt] (\x,5.9) -- (\x,6.5);
  \node[hw label, anchor=east] at (0.8,6.2) {帧到达};
  % 3 µs 里又到约 4 帧
  \draw[BookMuted, line width=0.6pt, <->] (1,5.3) -- (5.5,5.3);
  \node[hw label, anchor=south] at (3.25,5.4) {3 $\mu$s 里又到约 4 帧};
  % CPU：一次中断处理 3 µs，处理一帧又积下数帧
  \draw[BookAccent, fill=BookCodeBg, line width=0.8pt] (1,4.3) rectangle (5.5,4.9);
  \node at (3.25,4.6) {处理第 1 帧：3 $\mu$s};
  \draw[BookAccent, fill=BookCodeBg, line width=0.8pt] (5.5,4.3) rectangle (10,4.9);
  \node at (7.75,4.6) {处理第 2 帧：3 $\mu$s};
  \draw[BookAccent, fill=BookCodeBg, line width=0.8pt] (10,4.3) rectangle (13,4.9);
  \node[hw label] at (11.5,4.6) {处理第 3 帧 $\cdots$};
  \node[hw label, anchor=east] at (0.8,4.6) {中断处理};
  % 未处理积压：阶梯单调上涨（每个 3 µs 窗口净增约 3 帧）
  \draw[BookRed, line width=0.9pt] (1,1.2) -- (1,1.47) -- (2,1.47) -- (2,1.74) -- (3,1.74) -- (3,2.01) -- (4,2.01) -- (4,2.28) -- (5,2.28) -- (5,2.55) -- (5.5,2.55) -- (5.5,2.28) -- (6,2.28) -- (6,2.55) -- (7,2.55) -- (7,2.82) -- (8,2.82) -- (8,3.09) -- (9,3.09) -- (9,3.36) -- (10,3.36) -- (10,3.09) -- (11,3.09) -- (11,3.36) -- (12,3.36) -- (12,3.63) -- (13,3.63);
  \node[hw label, anchor=east] at (0.8,2.4) {未处理积压};
  \node[hw label] at (6.6,0.55) {帧每 0.67 $\mu$s 一帧，处理一帧又净增约 3 帧：越忙越处理不完};
\end{waveform}
\diagramnote{中断风暴（receive livelock）的时序（1 GbE 小帧线速）。帧每 0.67 $\mu$s 到达一帧，一次完整的中断处理——保存现场、读状态、递交协议栈、恢复现场——按保守的 3 $\mu$s 计：处理一帧的功夫线上又到约四帧，未处理积压单调上涨，需要的算力约 $1.49\times10^{6}\times3\ \mu$s 即 4.5 颗核的全部时间，协议栈和应用一行代码都轮不到。越忙越处理不完、吞吐反而崩；出路是把到达率除以批大小，即下一段的中断聚合。}

NAPI 把"中断驱动"改成"中断提示加轮询收尾"：第一帧触发中断后
关闭该队列的中断，驱动进入轮询，在预算内尽量多地收完积压的帧，
收干了再重新开中断——忙时纯轮询，闲时纯中断，两种状态自动切
换。网卡硬件侧的中断合并（interrupt moderation）提供两个阈值：
攒够 N 帧发一次中断，或距上次中断满 T µs 发一次，先到为准；低负
载时靠超时保证单帧不被无限积压。

补一个实现细节：NAPI 的轮询收包运行在软中断（softirq）上下文，
协议栈的 IP、TCP 处理大多也在其中。top 里看到的 si 一项的占比，
就是这条收包路径的 CPU 开销；它高到挤压应用，说明帧处理本身—
—而不只是中断进出——已经成为瓶颈，下一步手段是多队列分流，见
下一段。

权衡可以量化。设合并参数 N=32、T=40µs：小帧线速下 32 帧只需
$32\times0.672\approx21.5\mu s$，先于超时触发，中断率从约 149 万
次/秒降到 $1.49\times10^{6}/32\approx4.7$ 万次/秒，CPU 开销
$4.7\times10^{4}\times3\mu s\approx0.14$ 颗核——从中断风暴的 4.5
颗核降到约七分之一颗。代价是延迟：一帧平均多等约半批的时间，十
微秒量级。吞吐型业务把 N、T 调大；低延迟业务（行情、内存数据
库）调小甚至关闭合并，改用绑核加 busy-poll——拿一颗核 100\% 的
占用，换确定的几微秒时延。两个方向都对，区别只在业务把哪一头算
作成本。

把两种策略各画一条时序，这笔交换的形状如下：

\begin{waveform}[x=0.8cm,y=0.9cm]
  % 参考虚线：第 32 帧到达触发 / 等满 T 超时触发
  \draw[BookRule, line width=0.45pt, dashed] (8.75,1.35) -- (8.75,3.9);
  \draw[BookRule, line width=0.45pt, dashed] (12.0,1.35) -- (12.0,3.9);
  % 帧到达（线速，不聚合场景）
  \foreach \x in {1.0,1.25,...,8.9} \draw[BookTeal, line width=0.9pt] (\x,7.0) -- (\x,7.55);
  \node[hw label, anchor=east] at (0.8,7.28) {帧到达（线速）};
  % 中断·不聚合：每帧一次
  \draw[BookAccent, line width=0.9pt] (1.0,5.4) -- (9.0,5.4);
  \foreach \x in {1.0,1.25,...,8.9} \draw[BookAccent, line width=0.9pt] (\x,5.4) -- (\x,6.0) -- ({\x+0.13},6.0) -- ({\x+0.13},5.4);
  \node[hw label, anchor=east] at (0.8,5.7) {中断·不聚合};
  \node[hw label] at (6.4,4.9) {不聚合：每帧一次，$\approx$149 万次/秒，CPU $\approx$4.5 核};
  % 帧到达（聚合场景：线速 32 帧 + 低负载单帧）
  \foreach \x in {1.0,1.25,...,8.9} \draw[BookTeal, line width=0.9pt] (\x,3.2) -- (\x,3.75);
  \draw[BookTeal, line width=0.9pt] (10.2,3.2) -- (10.2,3.75);
  \node[hw label, anchor=east] at (0.8,3.48) {帧到达};
  \draw[BookMuted, line width=0.6pt, <->] (1.0,4.05) -- (8.75,4.05);
  \node[hw label] at (4.9,4.34) {攒满 32 帧 $\approx$21.5 $\mu$s（先于超时）};
  \draw[BookMuted, line width=0.6pt, <->] (10.2,4.05) -- (12.0,4.05);
  \node[hw label] at (11.1,4.34) {等满 $T$=40 $\mu$s};
  \node[hw label] at (10.4,2.92) {低负载：单帧到达};
  % 中断·聚合：攒满一批或超时，才发一次
  \draw[BookAccent, line width=0.9pt] (1.0,1.6) -- (8.75,1.6) -- (8.75,2.2) -- (9.05,2.2) -- (9.05,1.6) -- (12.0,1.6) -- (12.0,2.2) -- (12.3,2.2) -- (12.3,1.6) -- (12.6,1.6);
  \node[hw label, anchor=east] at (0.8,1.9) {中断·聚合};
  \node[hw label] at (6.3,0.95) {聚合 $N$=32 / $T$=40 $\mu$s：$\approx$4.7 万次/秒，CPU $\approx$0.14 核};
\end{waveform}
\diagramnote{中断聚合的时序对比。上方不聚合：每帧一次中断，1.49 Mpps 折合每秒约 149 万次中断进出，约 4.5 颗核的算力全花在保存与恢复现场上。下方聚合：攒满 32 帧（线速下约 21.5 $\mu$s，先于超时）或等满 T=40 $\mu$s——先到为准——才发一次中断，中断率降到约 4.7 万次/秒，CPU 开销约 0.14 核；代价是一帧平均多等约半个聚合窗口。}

\begin{longtable}{L{0.16\linewidth}L{0.32\linewidth}L{0.20\linewidth}L{0.20\linewidth}}
\toprule
参数 & 机制 & 偏向吞吐 & 偏向时延 \\
\midrule
帧数阈值 N & 攒够 N 帧发一次中断 & 调大（如 32--64） & 调小或置 1 \\
超时 T & 距上次中断满 T µs 即发 & 调大（几十 µs） & 调小或置 0 \\
轮询预算 & NAPI 每次最多收多少帧 & 调大 & 适中，防占核过久 \\
busy-poll & 应用自旋收帧，跳过中断 & 浪费核 & 时延最低且稳 \\
\bottomrule
\end{longtable}

观察手段是现成的：每秒读一次 /proc/interrupts 里各网卡队列的中
断计数，与流量对比就知道合并参数是否合身——中断率贴近 pps，说
明合并未生效；中断率过低而时延投诉上升，说明 N、T 过了头。
ethtool -c 与 -C 查询、调整这些参数。调网卡的第一步通常不是换硬
件，而是把中断率、环深、pps 三个数摆在一起看；三者相符，瓶颈才
轮到协议栈和应用。

\begin{lstlisting}
$ grep eth0 /proc/interrupts
 28:  1204431        0        0        0  IO-APIC  28-fasteoi  eth0-rx-0
 29:        0  1188204        0        0  IO-APIC  29-fasteoi  eth0-rx-1
 30:        0        0  1210022        0  IO-APIC  30-fasteoi  eth0-rx-2
 31:        0        0        0  1195330  IO-APIC  31-fasteoi  eth0-rx-3
# per-queue counters spread over 4 CPUs: RSS and affinity agree
\end{lstlisting}

读法：每个队列的中断计数应大致均匀地散在各核上；挤在一列说明
RSS 或绑核没有配对，整行全为 0 的队列说明中断亲和漏配。这个文
件也是下一话题的入口：多队列不是插上就有，是用对才有。

\topic{多队列与 RSS}单队列的天花板先算清楚：一个环只能由一颗核
服务——保序与 cache 局部性都要求如此——单核跑协议栈每秒约
1--2M 帧已是优化后的水平；10GbE 小帧是 14.88 Mpps，差一个数量
级。单队列网卡插上万兆链路，瓶颈不在线，在环。用第十章 Little
定律的读法：收帧路径的吞吐等于并行消费者数乘单消费者速率，在途
帧排得再多，消费者只有一颗核。

多队列网卡把 RX/TX 环做成很多对，常见 16--128 对，每对绑一个
MSI-X 中断向量、绑一颗核。接收侧的分流由硬件完成：RSS 对每帧的
五元组（源 IP、目的 IP、源端口、目的端口、协议号）计算 Toeplitz
哈希，取哈希低位查间接表，把帧定向到对应队列。同一条流哈希相
同，永远进同一队列，天然保序；不同的流散到不同核，并行处理。发
送侧对称：软件按流或按 CPU 选发送队列，各核各挂各的环，提交路
径无锁。

分流的全过程画成一张图，就是哈希定向加按核绑环：

\begin{pdfdiagram}
  \node[hw label, anchor=south] at (5.5,0.68) {同一条流哈希相同，永远进同一队列——天然保序};
  \node[hw state, minimum width=2.4cm] (nic) at (1.5,-1.28) {网卡收包\\10GbE 小帧\\14.88 Mpps};
  \node[hw compute, minimum width=2.8cm] (hash) at (4.35,-1.28) {RSS：五元组\\Toeplitz 哈希\\低位查间接表};
  \draw[hw bus] (nic.east) -- (hash.west);
  % 队列（示意 4 行，常见 16--128 对）
  \node[hw memory, minimum width=1.6cm, minimum height=0.55cm, inner sep=1pt] (q0) at (7.05,0) {RX 环 0};
  \node[hw memory, minimum width=1.6cm, minimum height=0.55cm, inner sep=1pt] (q1) at (7.05,-0.85) {RX 环 1};
  \node[hw label] at (7.05,-1.7) {$\vdots$};
  \node[hw memory, minimum width=1.6cm, minimum height=0.55cm, inner sep=1pt] (q15) at (7.05,-2.55) {RX 环 15};
  % 核
  \node[hw compute, minimum width=1.3cm, minimum height=0.55cm, inner sep=1pt] (c0) at (9.6,0) {核 0};
  \node[hw compute, minimum width=1.3cm, minimum height=0.55cm, inner sep=1pt] (c1) at (9.6,-0.85) {核 1};
  \node[hw label] at (9.6,-1.7) {$\vdots$};
  \node[hw compute, minimum width=1.3cm, minimum height=0.55cm, inner sep=1pt] (c15) at (9.6,-2.55) {核 15};
  % 哈希到队列：共享竖直干线，分层折出
  \draw[hw bus] (hash.east) -- (6.05,-1.28);
  \draw[BookMuted, line width=0.62pt] (6.05,0) -- (6.05,-2.55);
  \draw[hw bus] (6.05,0) -- (q0.west);
  \draw[hw bus] (6.05,-0.85) -- (q1.west);
  \draw[hw bus] (6.05,-2.55) -- (q15.west);
  % 队列到核：一对一
  \draw[hw bus] (q0.east) -- (c0.west);
  \draw[hw bus] (q1.east) -- (c1.west);
  \draw[hw bus] (q15.east) -- (c15.west);
  \node[hw label] at (5.5,-3.4) {每对环绑一个 MSI-X 向量、绑一颗核；单核约 1.5 Mpps，16 颗核并行才接得住 14.88 Mpps};
\end{pdfdiagram}
\diagramnote{RSS 接收分流：网卡对每帧的五元组算 Toeplitz 哈希，取低位查间接表，把帧定向到对应 RX 环；每对环绑一个 MSI-X 中断向量、绑一颗核，各核独立处理自己环上的帧，提交路径无锁。同一条流哈希相同、永远进同一队列，保序由哈希天然保证；并行度由队列对数决定——10GbE 小帧 14.88 Mpps 对单核约 1.5 Mpps，16 队列绑 16 核才接得住。}

配置按一遍数：10GbE 小帧 14.88 Mpps，单核约 1.5 Mpps，需要约 10
颗核起步，留余量配 16 队列、16 个 MSI-X 向量、绑 16 颗核。三个
配置必须一起看：队列数、间接表、中断亲和——队列开了 32 条而中
断全绑在核 0，等于单队列；irqbalance 之类的自动均衡在万兆以上
往往帮倒忙，绑核要按队列手工固定。

同样的算法推到更高速率：25GbE 小帧约 37.2 Mpps，40GbE 约 59.5
Mpps，100GbE 约 148.8 Mpps——单核约 1.5 Mpps 的处理率面前，需
要的核数按比例涨到约 25、40、100 颗，已经没有机器装得下。速率
每翻一番，"每帧固定成本"问题就尖锐一分：100GbE 时代还按"中断
加协议栈逐帧处理"的模型工作是不成立的，内核旁路与硬件卸载从可
选项变成必选项，这正是零拷贝一段要收尾的方向。

\begin{longtable}{L{0.18\linewidth}L{0.36\linewidth}L{0.34\linewidth}}
\toprule
组件 & 作用 & 配置不当的典型症状 \\
\midrule
队列对 & 收发各一条环，独立 DMA 与中断 & 队列少于核数，部分核闲置 \\
MSI-X 向量 & 每队列一个中断号 & 向量不足，多队列共享中断互拖 \\
RSS 间接表 & 哈希值到队列的映射 & 映射不均，各核负载歪斜 \\
中断亲和 & 把向量绑到指定核 & 全绑核 0，等于退回单队列 \\
\bottomrule
\end{longtable}

RSS 也有边界。它只认五元组：GRE、VxLAN 等隧道的外层包头都相
同，所有隧道流量挤进一条队列，较新的网卡支持按内层字段哈希来解
决。单条大象流永远只能跑一颗核的速度，这是保序的代价，只能靠多
开几条流来绕。与 RSS 同族的卸载还有 TSO 与 GRO：发送侧由网卡把
大块数据切成 MTU 帧，接收侧把同流相邻帧拼成大块再上交——它们
不改队列结构，改的是"每帧固定成本"里的帧数，与巨型帧是同一思
想的软硬件两种实现。

\topic{零拷贝的基本思想}先算拷贝一次的成本。传统收包路径上，网
卡 DMA 把帧写进内核缓冲区，这一步不占 CPU；协议栈处理后再由
copy\_to\_user 把数据复制一份到应用缓冲区——这一次是 CPU 逐字
节的复制，读一遍、写一遍。10Gb/s 的流量约 1.25GB/s，一次拷贝产
生 2.5GB/s 的内存流量；收发各拷一次就是 5GB/s，约占一条
DDR4-3200 通道（25.6GB/s）可用带宽的两成，还没算 cache 被冲刷的
间接代价。用第十章 roofline 的话说：每次无谓的拷贝都在消耗全机
共享的带宽屋顶，留给应用的屋顶相应变矮。

零拷贝的思路是让数据在内存里只躺一份，流转的是"这块缓冲区归
谁"。sendfile 让文件页缓存直接对接 socket，配合支持
scatter-gather 的网卡，连"协议头加数据"的拼接都由网卡 gather
DMA 完成，CPU 一个字节不搬；splice 在两个文件描述符之间搬运页
指针；mmap 让应用直接读写内核页，省掉复制但引入页错误与同步的
成本；io\_uring 的固定缓冲区与 AF\_XDP 更进一步，把缓冲区的所有
权约定直接开放给用户态。注意"零拷贝"省的从来不是 DMA 那一次
——设备搬运本来就绕不开——省的是内核与用户态之间那一次 CPU
复制。

\begin{lstlisting}
#include <sys/sendfile.h>

/* stream a file to a socket without copying through user space */
ssize_t sendfile(int out_fd, int in_fd, off_t *offset, size_t count);

/* kernel path: page cache pages -> socket buffer -> NIC gather DMA;
   file contents never pass through a CPU register */
\end{lstlisting}

两条路径并排画出来，零拷贝省掉的是哪一段就一目了然：

\begin{pdfdiagram}
  % 传统路径：两次 CPU 拷贝
  \node[hw label, anchor=west] at (0.4,0.98) {传统：页缓存 $\to$ 应用缓冲 $\to$ socket 缓冲，两次 CPU 拷贝};
  \node[hw memory, minimum width=2.0cm, minimum height=0.8cm] (pc1) at (1.4,0) {页缓存};
  \node[hw memory, minimum width=2.0cm, minimum height=0.8cm] (ub) at (4.2,0) {应用缓冲};
  \node[hw memory, minimum width=2.0cm, minimum height=0.8cm] (sb1) at (7.0,0) {socket 缓冲};
  \node[hw memory, minimum width=2.0cm, minimum height=0.8cm] (nic1) at (9.8,0) {网卡};
  \draw[hw bus] (pc1.east) -- (ub.west);
  \draw[hw bus] (ub.east) -- (sb1.west);
  \draw[hw bus] (sb1.east) -- (nic1.west);
  \node[hw label, anchor=south] at (2.8,0.62) {① CPU 拷贝};
  \node[hw label, anchor=south] at (5.6,0.62) {② CPU 拷贝};
  \node[hw label, anchor=south] at (8.4,0.62) {③ DMA};
  % 零拷贝路径
  \node[hw label, anchor=west] at (0.4,-0.98) {零拷贝 sendfile：流转的是缓冲区归属，CPU 不搬数据};
  \node[hw memory, minimum width=2.0cm, minimum height=0.8cm] (pc2) at (1.4,-1.95) {页缓存};
  \node[hw memory, minimum width=2.6cm, minimum height=0.8cm] (sb2) at (5.6,-1.95) {socket 缓冲\\只记指针};
  \node[hw memory, minimum width=2.0cm, minimum height=0.8cm] (nic2) at (9.8,-1.95) {网卡};
  \draw[hw return] (pc2.east) -- (sb2.west);
  \draw[hw bus] (sb2.east) -- (nic2.west);
  \node[hw label, anchor=south] at (2.85,-1.33) {只改归属};
  \node[hw label, anchor=south] at (8.35,-1.33) {gather DMA};
  \node[hw label] at (5.5,-2.9) {10 Gb/s 下两次 CPU 拷贝产生 5 GB/s 内存流量，约占一条 DDR4-3200 通道（25.6 GB/s）的两成};
\end{pdfdiagram}
\diagramnote{发送一个文件的两种路径。传统路径上数据两次经过 CPU：页缓存拷进应用缓冲、应用缓冲再拷进 socket 缓冲，每次都是读一遍写一遍，10 Gb/s 的流量就此产生 5 GB/s 的内存流量。sendfile 让页缓存直接对接 socket，缓冲区只改归属、内容原地不动，网卡用 gather DMA 把协议头与数据拼接发出，CPU 一个字节不搬——省的是内核与用户态之间那次复制，设备 DMA 那一步本来就有。}

\begin{longtable}{L{0.18\linewidth}L{0.34\linewidth}L{0.34\linewidth}}
\toprule
机制 & 省掉的是哪次拷贝 & 适用场景 \\
\midrule
sendfile & 文件页缓存到用户态的往返 & 静态文件服务、CDN \\
splice & 任意两个 fd 之间的搬运 & 代理、转发管道 \\
mmap & 读方向的复制 & 随机访问大文件 \\
io\_uring 固定缓冲 & 每次 I/O 的缓冲区注册与复制 & 高频异步 I/O \\
AF\_XDP / DPDK & 整个内核协议栈路径 & 极端 pps 的专用应用 \\
\bottomrule
\end{longtable}

工程含义与边界。零拷贝的本质是所有权转移代替内容复制，正是第五
章"缓冲区何时归 CPU、何时归设备"那条问题在协议栈内部的延长
线。它不是免费午餐：页共享带来生命周期管理与写时复制的复杂度，
数据量小时固定开销可能盖过节省；它适用的条件是大块数据、高带
宽、CPU 比带宽更紧张——这正是存储服务器、CDN、反向代理人手一
份 sendfile 的原因。从本章视角看，零拷贝与中断聚合、RSS 是同一
族手段：网络栈的性能工作，大半是在给"每帧固定成本"与"每字节
搬运成本"这两项开销找出路。

旁路也不是没有代价。把协议栈整个绕开，意味着 TCP 的拥塞控制、
分片、校验和、防火墙一并让位，应用要么自己实现，要么声明不需
要；内核协议栈占的 CPU，换来的正是这些不免费的功能。工程上成熟
的做法是分层：极端 pps 的入口（行情网关、负载均衡）用旁路或
AF\_XDP，业务主体仍走内核栈加零拷贝。第五章那句提醒在这里同样
成立：约定一旦绕开，原来由内核代管的所有权、可见性与错误处理，
全部换成应用自己的功课。

\section{五、外设队列综合案例与尾延迟}

前四节分别看了三类设备的介质与各自的队列结构，本节把它们汇合到一个
问题上：队列机制怎样决定外设实际表现出来的性能。

\topic{NVMe 的队列模型}
第五章讲描述符环时给过一般结构：环不动、指针动，doorbell 只通知、不
携带数据，完成与否以设备写回的状态为准。NVMe 把这套结构定型为规范，
并针对多核做了两处关键改动。第一处改动是队列成对：NVMe 的基本单元
是一对队列——提交队列 SQ（submission queue）与完成队列 CQ
（completion queue）。SQ 项是 64 字节的命令：操作码（读、写、
flush）、起始 LBA、块数，以及 PRP 数据指针列表，逐页给出数据在主机
内存里的物理地址；CQ 项是 16 字节的完成记录：命令 ID、状态码，以及
一个相位位——队列回绕一圈相位翻一次，主机据此区分新旧完成项，不必
比较指针。提交与完成分成两个环，各走各的指针：设备消费 SQ 不必等主
机回收 SQ，完成信息也不挤在命令项里写回。这是比第五章通用环更彻底
的分工。

一次提交的完整序列有七步，逐步对应第五章的每一条约定：

\begin{lstlisting}
1. 驱动把 64 B 命令项写入本核 SQ 的尾槽
2. 屏障：保证命令项先于通知对设备可见
3. 写 SQ 尾 doorbell（一次 MMIO 写，内容只是新尾指针）
4. 设备看到指针差，DMA 取走新命令项
5. 设备执行：查 FTL、读 NAND，按 PRP 地址把数据 DMA 进内存
6. 设备写 16 B CQ 完成项，相位位翻转
7. 设备发 MSI-X 中断；驱动读 CQ、按命令 ID 找回请求并回收，
   最后写 CQ 头 doorbell，把完成槽位还给设备
\end{lstlisting}

第 3 步的 doorbell 仍然不携带数据：设备靠首尾指针的差补读新命令，偶
尔错过一次 doorbell 并无妨碍；第 6、7 步的次序也不许交换：先写完成
项、再发中断，主机在中断处理里读到的一定是已经写好的 CQ 项。第五章
"描述符先可见、再通知"和"完成以写回状态为准"这两条规则，在这里原
样成立，只是被写进了设备规范，由硬件而不是驱动员的细心来保证。

把这条链路画成图，主机内存、doorbell 寄存器与设备之间的
分工如下：

\begin{pdfdiagram}
  \node[hw state, minimum width=1.9cm] (cpu) at (0.9,0.55) {CPU\\（本核驱动）};
  % 主机内存：SQ 环与 CQ 环
  \draw[BookAmber, line width=0.42pt, rounded corners=2pt] (2.6,-1.3) rectangle (7.2,1.35);
  \node[hw label, anchor=west] at (2.75,1.08) {SQ：64 B 命令项环};
  \node[hw label, anchor=east] at (7.1,0.12) {主机内存（DRAM）};
  \node[hw state, minimum width=0.78cm, minimum height=0.5cm] (sq1) at (3.35,0.55) {命令};
  \node[hw state, minimum width=0.78cm, minimum height=0.5cm] at (4.35,0.55) {命令};
  \node[hw state, minimum width=0.78cm, minimum height=0.5cm] at (5.35,0.55) {命令};
  \node[hw state, minimum width=0.78cm, minimum height=0.5cm] (sq4) at (6.35,0.55) {$\cdots$};
  \node[hw state, minimum width=0.78cm, minimum height=0.5cm] at (3.35,-0.55) {完成};
  \node[hw state, minimum width=0.78cm, minimum height=0.5cm] at (4.35,-0.55) {完成};
  \node[hw state, minimum width=0.78cm, minimum height=0.5cm] at (5.35,-0.55) {完成};
  \node[hw state, minimum width=0.78cm, minimum height=0.5cm] (cq4) at (6.35,-0.55) {$\cdots$};
  \node[hw label, anchor=west] at (2.75,-1.05) {CQ：16 B 完成项环};
  \node[hw control, minimum width=1.6cm] (db) at (8.3,0.45) {doorbell\\寄存器};
  \node[hw compute, minimum width=1.9cm] (dev) at (10.45,0.45) {NVMe 控制器\\FTL + DMA};
  % ① 驱动把命令项写进本核 SQ 尾槽
  \draw[hw bus] (cpu.east) -- (sq1.west);
  \node[hw label, anchor=west] at (2.75,0.12) {① 写命令项};
  % ② 屏障后写 SQ 尾 doorbell（MMIO，控制）
  \draw[hw return] (cpu.south) -- (0.9,-2.0) -- (8.3,-2.0) -- (db.south);
  \node[hw label, above] at (4.5,-2.0) {② 写 SQ 尾 doorbell（MMIO）};
  % ③ 设备按指针差 DMA 取走命令
  \draw[hw bus] (sq4.north) -- (6.35,1.85) -- (10.45,1.85) -- (dev.north);
  \node[hw label, above] at (8.3,1.85) {③ DMA 取命令};
  % doorbell 通知设备（只携带新尾指针）
  \draw[hw return] (db.east) -- (dev.west);
  % ④ 设备写完成项进 CQ（device.south 竖直干线共享）
  \draw[hw bus] (dev.south) -- (10.45,-1.5) -- (6.35,-1.5) -- (cq4.south);
  \node[hw label, below] at (7.3,-1.5) {④ 写完成项};
  % ⑤ MSI 中断回本核
  \draw[hw return] (dev.south) -- (10.45,-2.6) -- (0.3,-2.6) -- (0.3,0.07);
  \node[hw label, below] at (5.0,-2.6) {⑤ MSI 中断回本核};
\end{pdfdiagram}
\diagramnote{NVMe 的提交—完成链路。驱动把 64 B 命令项写进本核 SQ 尾槽（①），屏障之后写 SQ 尾 doorbell（②）——doorbell 只携带新尾指针；设备按指针差 DMA 取走命令（③），执行后把 16 B 完成项写进 CQ（④），此时数据已按 PRP 地址 DMA 进入内存，最后以 MSI-X 中断通知本核（⑤）。SQ 与 CQ 各自独立、每核一对，提交路径上没有需要上锁的共享状态。}

第二处改动是多队列。规范允许最多 65535 对 I/O 队列，每对深度上限
65535；真实盘实现的对数远小于此（几十到几千），但已足够给每个 CPU
核配一对。每核一对的用意是免锁：各核写自己的 SQ、收自己的 CQ，尾指
针是核私有的，提交路径上没有任何共享状态需要上锁。这不是锁实现得好，
而是根本没有共享——第十章谈多核时说过，把共享拆开比把锁做快更值钱，
NVMe 的多队列是这句话在外设接口上的直接应用。配合 MSI-X 的多向量中
断，完成也能按队列分发回本核处理，提交与完成两条路径都不出核。

把每核一对的结构画出来，免锁与不出核是同一件事的两个侧面：

\begin{pdfdiagram}
  % 左：各核
  \node[hw compute, minimum width=1.1cm, minimum height=0.62cm] (k0) at (0.85,0.5) {核 0};
  \node[hw compute, minimum width=1.1cm, minimum height=0.62cm] (k1) at (0.85,-0.6) {核 1};
  \node[hw label] at (0.85,-1.55) {$\vdots$};
  \node[hw compute, minimum width=1.1cm, minimum height=0.62cm] (kn) at (0.85,-2.5) {核 $n$};
  % 中：每核一对 SQ/CQ（主机内存）
  \node[hw memory, minimum width=3.1cm] (p0) at (4.05,0.5) {SQ 0：命令环\\CQ 0：完成环};
  \node[hw memory, minimum width=3.1cm] (p1) at (4.05,-0.6) {SQ 1：命令环\\CQ 1：完成环};
  \node[hw label] at (4.05,-1.55) {$\vdots$};
  \node[hw memory, minimum width=3.1cm] (pn) at (4.05,-2.5) {SQ $n$：命令环\\CQ $n$：完成环};
  \node[hw compute, minimum width=2.6cm] (ctl) at (8.35,-1.0) {NVMe 控制器\\FTL＋DMA 引擎};
  % 提交：命令项 + doorbell
  \draw[hw bus] (k0.east) -- (p0.west);
  \draw[hw bus] (k1.east) -- (p1.west);
  \draw[hw bus] (kn.east) -- (pn.west);
  \node[hw label, anchor=south] at (1.95,1.0) {写命令项};
  \draw[BookMuted, line width=0.62pt] (6.15,0.5) -- (6.15,-2.5);
  \draw[hw bus] (p0.east) -- (6.15,0.5);
  \draw[hw bus] (p1.east) -- (6.15,-0.6);
  \draw[hw bus] (pn.east) -- (6.15,-2.5);
  \draw[hw bus] (6.15,-1.0) -- (ctl.west);
  \node[hw label, anchor=east] at (6.0,-1.15) {DMA 取命令};
  % doorbell：只带新尾指针
  \draw[hw return] (kn.south) -- (0.85,-3.35) -- (8.35,-3.35) -- (ctl.south);
  \node[hw label, anchor=north] at (4.6,-3.42) {doorbell：只带新尾指针（每队列一对寄存器）};
  % 完成：写 CQ + MSI-X 回本核
  \draw[hw return] (ctl.north) -- (8.35,1.35) -- (4.05,1.35) -- (p0.north);
  \node[hw label, anchor=south] at (6.2,1.42) {写 CQ 完成项，MSI-X 回本核};
  \node[hw label] at (4.6,-4.25) {尾指针核私有：提交路径上没有共享状态；规范上限 65535 对、每对深 65535};
\end{pdfdiagram}
\diagramnote{NVMe 的多队列结构：每核一对 SQ/CQ，尾指针核私有，doorbell 只携带新尾指针；设备经 DMA 取走命令、写 CQ 完成项，再以 MSI-X 按队列分发回本核——提交与完成两条路径都不出核。免锁不是锁实现得好，而是提交路径上根本没有共享状态。规范允许最多 65535 对队列、每对深 65535，真实盘实现几十到几千对，常用深度 32 至 1024。}

对照 AHCI，就能看出这两处改动解决了什么。AHCI 整台控制器只有一张命
令表，32 个槽位（NCQ 深度 32），一个完成中断源：所有核的提交都抢同
一张表的锁，完成中断也集中在一处。在 SATA 的 600 MB/s、约十万 IOPS
的上限下，这张表不是瓶颈；PCIe SSD 把带宽抬高一个数量级、把 IOPS
抬高两个数量级，提交路径上的共享状态先成为瓶颈。NVMe 的多队列不是
锦上添花，而是把 AHCI 那张唯一的表拆开之后的必然形态。

把两代接口的提交结构并排画出来，瓶颈的出处一目了然：

\begin{pdfdiagram}
  % 分隔线
  \draw[BookRule, line width=0.45pt, dashed] (5.35,0.9) -- (5.35,-3.0);
  % 左：AHCI
  \node[hw label] at (2.55,0.95) {AHCI：整控制器一张命令表（32 槽）};
  \node[hw compute, minimum width=0.9cm, minimum height=0.55cm] (c0) at (0.55,0.35) {核 0};
  \node[hw compute, minimum width=0.9cm, minimum height=0.55cm] (c1) at (0.55,-0.35) {核 1};
  \node[hw compute, minimum width=0.9cm, minimum height=0.55cm] (c2) at (0.55,-1.05) {核 2};
  \node[hw compute, minimum width=0.9cm, minimum height=0.55cm] (c3) at (0.55,-1.75) {核 3};
  \draw[BookMuted, line width=0.62pt] (1.35,0.35) -- (1.35,-1.75);
  \draw[hw bus] (c0.east) -- (1.35,0.35);
  \draw[hw bus] (c1.east) -- (1.35,-0.35);
  \draw[hw bus] (c2.east) -- (1.35,-1.05);
  \draw[hw bus] (c3.east) -- (1.35,-1.75);
  \node[hw control, minimum width=1.3cm] (lk) at (2.35,-0.7) {表锁};
  \node[hw memory, minimum width=2.0cm] (tb) at (4.1,-0.7) {命令表\\32 槽（NCQ）};
  \draw[hw bus] (1.35,-0.7) -- (lk.west);
  \draw[hw bus] (lk.east) -- (tb.west);
  \node[hw label] at (2.55,-2.6) {多核提交抢同一把表锁\\完成中断集中一处};
  % 右：NVMe
  \node[hw label] at (8.6,0.95) {NVMe：每核一对 SQ/CQ};
  \node[hw compute, minimum width=0.9cm, minimum height=0.55cm] (d0) at (6.25,0.35) {核 0};
  \node[hw compute, minimum width=0.9cm, minimum height=0.55cm] (d1) at (6.25,-0.35) {核 1};
  \node[hw compute, minimum width=0.9cm, minimum height=0.55cm] (d2) at (6.25,-1.05) {核 2};
  \node[hw compute, minimum width=0.9cm, minimum height=0.55cm] (d3) at (6.25,-1.75) {核 3};
  \node[hw memory, minimum width=1.9cm, minimum height=0.55cm] (q0) at (8.05,0.35) {SQ/CQ 0};
  \node[hw memory, minimum width=1.9cm, minimum height=0.55cm] (q1) at (8.05,-0.35) {SQ/CQ 1};
  \node[hw memory, minimum width=1.9cm, minimum height=0.55cm] (q2) at (8.05,-1.05) {SQ/CQ 2};
  \node[hw memory, minimum width=1.9cm, minimum height=0.55cm] (q3) at (8.05,-1.75) {SQ/CQ 3};
  \draw[hw bus] (d0.east) -- (q0.west);
  \draw[hw bus] (d1.east) -- (q1.west);
  \draw[hw bus] (d2.east) -- (q2.west);
  \draw[hw bus] (d3.east) -- (q3.west);
  \draw[BookMuted, line width=0.62pt] (9.2,0.35) -- (9.2,-1.75);
  \draw[hw bus] (q0.east) -- (9.2,0.35);
  \draw[hw bus] (q1.east) -- (9.2,-0.35);
  \draw[hw bus] (q2.east) -- (9.2,-1.05);
  \draw[hw bus] (q3.east) -- (9.2,-1.75);
  \node[hw compute, minimum width=1.7cm] (nc) at (10.25,-0.7) {NVMe\\控制器};
  \draw[hw bus] (9.2,-0.7) -- (nc.west);
  \node[hw label] at (8.6,-2.6) {核私有指针，提交免锁\\MSI-X 按核分发完成};
\end{pdfdiagram}
\diagramnote{AHCI 与 NVMe 的提交结构对照。AHCI 整台控制器只有一张 32 槽命令表（NCQ 深度 32）：所有核的提交抢同一把表锁，完成中断也集中在一处；SATA 的 600 MB/s、约十万 IOPS 下这张表不是瓶颈，PCIe SSD 把带宽抬高一个数量级、IOPS 抬高两个数量级之后，提交路径上的共享状态先成为瓶颈。NVMe 把表拆开：每核一对 SQ/CQ，尾指针核私有、提交免锁，完成由 MSI-X 按核分发——多队列是把那张唯一的表拆开之后的必然形态。}

\begin{longtable}{L{0.16\linewidth}L{0.30\linewidth}L{0.34\linewidth}}
\toprule
维度 & AHCI（SATA 时代） & NVMe（PCIe 时代） \\
\midrule
队列数 & 整控制器一张命令表 & 规范上限 65535 对 SQ/CQ \\
每队深度 & 32（NCQ 槽位） & 上限 65535，常用 32 至 1024 \\
提交路径 & 共享表，多核提交要抢锁 & 每核一对，核私有指针，免锁 \\
完成通知 & 单一中断源 & MSI-X 多向量，可按核分发 \\
典型上限 & 600 MB/s、约十万 IOPS & 数 GB/s、几十万到百万 IOPS \\
\bottomrule
\end{longtable}

深度上限 65535 也不意味着要用满。队列深度的用途是维持在途请求数，
第十章的 Little 定律给出该维持多少；用满上限买到的是排队而不是吞
吐——本节第四主题会把这笔账算到具体数字。

\topic{案例：一次 NVMe 读的完整 trace}
把一个进程执行 \code{read(fd, buf, 4096)}、且数据不在页缓存的全过
程按时序列出。设备取一块典型的 PCIe 3.0 x4 NVMe SSD，请求是 4 KB
随机读。

\begin{longtable}{L{0.06\linewidth}L{0.44\linewidth}L{0.14\linewidth}L{0.16\linewidth}}
\toprule
步 & 动作 & 耗时量级 & 执行者 \\
\midrule
1 & 系统调用陷入内核，检查参数与文件描述符 & 约 1 $\mu$s & CPU \\
2 & VFS 查页缓存：未命中，分配空闲页框 & 微秒级 & CPU \\
3 & 文件系统把文件偏移翻译成设备 LBA & 微秒级 & CPU \\
4 & 块层生成请求，调度器尝试与相邻请求合并 & 微秒级 & CPU \\
5 & 驱动把请求编成 64 B 命令项，写入本核 SQ 尾槽 & 微秒级 & CPU \\
6 & 屏障之后写 SQ 尾 doorbell（MMIO 写） & 百纳秒级 & CPU \\
7 & 设备 DMA 取走命令项 & 微秒级 & 设备 \\
8 & 查 FTL 映射，读 NAND 页，ECC 校验 & 约 80 $\mu$s & 介质 \\
9 & 数据按 PRP 地址 DMA 写入第 2 步的页框 & 约 1 $\mu$s & 设备 \\
10 & 设备写 16 B CQ 完成项，相位位翻转 & 微秒级 & 设备 \\
11 & 设备发 MSI-X，中断控制器转成 CPU 中断 & 微秒级 & 设备 \\
12 & 中断处理读 CQ，找回请求，标记完成，唤醒进程 & 微秒级 & CPU \\
13 & 页缓存数据拷入用户缓冲，\code{read} 返回 4096 & 微秒级 & CPU \\
\bottomrule
\end{longtable}

总时延约 100 $\mu$s，其中第 8 步介质本身占去八成，软件与队列合计约
两成——这就是 NVMe 时代的比例：协议和驱动的开销已被压到介质时延的
零头，继续优化软件路径收益有限，优化介质的访问形态（对齐、合并、足
够的在途深度）才是主战场。这张表也是排障地图：慢在哪个量级，就查哪
一列——微秒级的步骤出问题多半是驱动或内核配置，毫秒级的问题才轮到
介质与垃圾回收。

第 6 步值得停一下：先写命令项、再写 doorbell，中间隔一道屏障，这正
是第五章"描述符先可见、再通知"的约定；顺序反了，设备可能读到半写
入的命令项。第 11 步也值得一提：MSI 不是一根专用电平线，而是设备向
中断控制器的特定地址做一次写事务。第八章讨论过沿导线传播的信号要付
出走线与时序的代价；MSI 干脆把中断变成数据流里的一次写，省掉专用引
脚，还能在写数据里携带向量号，MSI-X 于是能按队列把中断分发到不同
核——中断从"一根线"演化成"一条消息"，和多队列是同一个设计方向。

同一件事在第六章的 BIOS 路径上长什么样，对照一下就清楚了：

\begin{longtable}{L{0.14\linewidth}L{0.34\linewidth}L{0.34\linewidth}}
\toprule
维度 & int 13h（第六章 BIOS 路径） & NVMe 路径 \\
\midrule
提交 & 参数装寄存器，执行 \code{int 0x13} & 命令项写内存队列，doorbell 通知 \\
地址 & CHS 柱面/磁头/扇区 & LBA 加 PRP 物理页列表 \\
等待 & BIOS 轮询状态寄存器，独占 CPU & 进程睡眠，完成由 MSI 唤醒 \\
并发 & 一次一个在途 & 几十上百在途，多队列并行 \\
搬运 & 软盘走 8237 DMA，硬盘 PIO 或总线主控 & 设备直接 DMA 到页框 \\
\bottomrule
\end{longtable}

骨架是同一个：准备命令、通知设备、DMA 搬数、确认完成。int 13h 是这
套骨架的同步简化版——提交即等待，全程没有队列，CPU 停在 BIOS 例程
里轮询；NVMe 是它的并发完整版——提交与完成解耦，等待换成中断，在
途请求从一个变成几百个。从 8237 的一个 DMA 通道到每核一对队列，外
设接口四十年的演进主线，就是把"一次一单"改成"持续事务流"，而描
述符与 DMA 这两个零件自始至终没有换过。

\topic{尾延迟：P99 从哪里来}
前面反复出现的"典型时延 100 $\mu$s"只是中位数附近的值。第十章已
经立下规矩：时延是一个分布，平均值会撒谎。外设是把这条规矩演绎得最
充分的地方——慢请求不是随机噪声，而是有明确物理来源的事件，每一类
都对应本章前几节讲过的机制：

\begin{longtable}{L{0.16\linewidth}L{0.36\linewidth}L{0.14\linewidth}L{0.20\linewidth}}
\toprule
长尾来源 & 物理机制 & 典型幅度 & 触发条件 \\
\midrule
GC 擦除 & 写请求排在一次块擦除与有效页搬运后面 & 毫秒级 & 空闲块耗尽、盘接近写满 \\
队列拥塞 & 深队列高负载下排队，第十章的 $1/(1-\rho)$ & 百微秒级 & 深度过大或负载接近饱和 \\
SMR 重写 & 区域写满触发整段重写（第三节） & 十毫秒级 & 随机写命中已写区域 \\
读重试 & 阈值电压漂移，ECC 纠不动，换挡重读 & 数十微秒至毫秒 & 老盘、高温、久置数据 \\
映射回读 & FTL 映射在 DRAM 中未命中，回读 NAND 映射页 & 百微秒级 & 冷地址、大跨度跳跃 \\
\bottomrule
\end{longtable}

这些事件不常发生，所以中位数很好看；但一旦发生就是十倍、百倍的慢。
消费级 NVMe SSD 上 4 KB 随机读的典型分布是：

\begin{longtable}{L{0.14\linewidth}L{0.16\linewidth}L{0.20\linewidth}L{0.32\linewidth}}
\toprule
百分位 & 时延 & 相对中位数 & 直观含义 \\
\midrule
P50 & 100 $\mu$s & 1 倍 & 典型请求 \\
P99 & 1 ms & 10 倍 & 每 100 个里约 1 个这么慢 \\
P999 & 10 ms & 100 倍 & 每 1000 个里约 1 个这么慢 \\
\bottomrule
\end{longtable}

把这张分布画成直方图，长尾的形状比数字更直接：

\begin{pdfdiagram}
  % 坐标轴
  \draw[BookMuted, line width=0.6pt, ->] (0.5,0) -- (11.0,0);
  \draw[BookMuted, line width=0.6pt, ->] (0.5,0) -- (0.5,4.0);
  \node[hw label, anchor=south] at (0.5,4.05) {请求数};
  \node[hw label] at (6.9,-0.3) {时延（横轴按对数间距）};
  % 直方图：主峰在 100 µs 附近，长尾拖向 1 ms、10 ms
  \filldraw[fill=BookTeal!25, draw=BookTeal, line width=0.5pt] (0.8,0) rectangle (1.3,0.7);
  \filldraw[fill=BookTeal!25, draw=BookTeal, line width=0.5pt] (1.35,0) rectangle (1.85,3.2);
  \filldraw[fill=BookTeal!25, draw=BookTeal, line width=0.5pt] (1.95,0) rectangle (2.45,1.5);
  \filldraw[fill=BookTeal!25, draw=BookTeal, line width=0.5pt] (2.55,0) rectangle (3.05,0.8);
  \filldraw[fill=BookTeal!25, draw=BookTeal, line width=0.5pt] (3.25,0) rectangle (3.75,0.5);
  \filldraw[fill=BookTeal!25, draw=BookTeal, line width=0.5pt] (4.05,0) rectangle (4.55,0.35);
  \filldraw[fill=BookTeal!25, draw=BookTeal, line width=0.5pt] (5.05,0) rectangle (5.55,0.25);
  \filldraw[fill=BookTeal!25, draw=BookTeal, line width=0.5pt] (6.25,0) rectangle (6.75,0.18);
  \filldraw[fill=BookTeal!25, draw=BookTeal, line width=0.5pt] (7.65,0) rectangle (8.15,0.13);
  \filldraw[fill=BookTeal!25, draw=BookTeal, line width=0.5pt] (9.15,0) rectangle (9.65,0.1);
  \filldraw[fill=BookTeal!25, draw=BookTeal, line width=0.5pt] (10.15,0) rectangle (10.65,0.07);
  % 百分位参考线与平均值
  \draw[BookMuted, line width=0.5pt, dashed] (1.6,0) -- (1.6,3.75);
  \node[hw label, anchor=west] at (1.72,3.9) {P50 = 100 $\mu$s};
  \draw[BookRed, line width=0.5pt, dashed] (2.5,0) -- (2.5,3.35);
  \node[hw label, anchor=west, text=BookRed] at (2.62,3.5) {平均值 $\approx$199 $\mu$s};
  \draw[BookRed, line width=0.6pt, ->] (1.72,2.9) -- (2.4,2.9);
  \node[hw label, anchor=west] at (2.62,2.9) {被长尾拉偏};
  \draw[BookMuted, line width=0.5pt, dashed] (9.4,0) -- (9.4,1.5);
  \node[hw label, anchor=south] at (9.4,1.55) {P99 = 1 ms};
  \draw[BookMuted, line width=0.5pt, dashed] (10.4,0) -- (10.4,1.1);
  \node[hw label, anchor=south] at (10.4,1.15) {P999 = 10 ms};
  % 横轴刻度与长尾来源
  \node[hw label, anchor=north] at (1.6,-0.12) {100 $\mu$s};
  \node[hw label, anchor=north] at (9.4,-0.12) {1 ms};
  \node[hw label, anchor=north] at (10.4,-0.12) {10 ms};
  \node[hw label] at (6.2,0.85) {长尾：GC 擦除、队列拥塞、SMR 重写、读重试};
\end{pdfdiagram}
\diagramnote{消费级 NVMe SSD 4 KB 随机读的时延分布（横轴按对数间距）。主峰在 P50 = 100 $\mu$s；GC 擦除、队列拥塞、SMR 重写、读重试这类不常发生、一发生就是十倍百倍慢的事件，把分布拖出长尾：P99 = 1 ms，P999 = 10 ms。平均值被长尾向右拉偏，离中位数越来越远——所以外设时延必须带百分位一起报。}

平均时延把这些全藏了起来。算一遍第十章的同款例子：100 个请求里 99
个花 100 $\mu$s、1 个花 10 ms，平均值是
$(99\times100+10000)/100=199\ \mu$s——平均看上去只比中位数差一倍，
可那一个请求慢了 100 倍。若服务目标是"99\% 的请求快于
200 $\mu$s"，一份只报平均值 199 $\mu$s 的报告，会把一次彻底的违约
藏进"接近达标"的表象里。用户记住的从来不是平均那次，而是最慢那
次：页面偶尔卡住三秒，原因往往就是排在 P999 上的那个请求撞上了 GC
擦除或 SMR 重写。

工程含义有三条。其一，报外设时延必须带百分位：平均值至少配 P99，交
互与实时负载再看 P999，只报平均值的测试报告不可信。其二，容量规划
按 P99 反推负载上限：让 P99 保持在目标内的负载，通常只有平均时延开
始明显变坏时的一半左右——这与第十章"稳态利用率压在五到七成"是同
一条纪律在外设上的形式。其三，长尾不能靠平均优化消掉，只能按来源逐
个治理：GC 尖峰靠预留空闲块和 TRIM，队列拥塞靠限深或限流，SMR 重写
靠改变写入形态——这正是本节后面两个主题的内容。

\topic{队列深度与吞吐时延的权衡}
第十章的 Little 定律 $L=\lambda W$ 用在外设上的形式是：在途请求数
等于 IOPS 乘单请求时延。介质时延 $W$ 由物理决定（4 KB 随机读约
100 $\mu$s），吞吐目标 $\lambda$ 由应用决定，队列深度是主机侧唯一
的旋钮。两个方向都有代价。

深度太小，喂不饱设备。QD=1 时 $\lambda=1/W=10^{4}$ IOPS，按 4 KB
折算只有 40 MB/s——一块标称 3.5 GB/s 的盘只用了约 1\%。深度加大，
吞吐近似线性上升，直到撞上设备自身的并行度上限（NAND 通道数与控制
器能力）。设某盘上限为 50 万 IOPS，拐点的位置由
$L=\lambda W=5\times10^{5}\times10^{-4}=50$ 唯一决定。越过拐点再加
深度，吞吐不变，多出来的在途全部变成排队：QD=256 时平均时延约
$256/(5\times10^{5})=512\ \mu$s，是介质时延的五倍多。

\begin{longtable}{L{0.12\linewidth}L{0.18\linewidth}L{0.18\linewidth}L{0.34\linewidth}}
\toprule
队列深度 & 可达 IOPS & 平均时延 & 位置 \\
\midrule
1 & 1 万 & 100 $\mu$s & 深度限制，带宽饿着 \\
4 & 4 万 & 100 $\mu$s & 深度限制 \\
16 & 16 万 & 100 $\mu$s & 深度限制 \\
32 & 32 万 & 100 $\mu$s & 接近拐点 \\
64 & 50 万（封顶） & 约 128 $\mu$s & 刚过拐点，开始排队 \\
128 & 50 万 & 约 256 $\mu$s & 排队区，吞吐不再涨 \\
256 & 50 万 & 约 512 $\mu$s & 排队区，纯买时延 \\
\bottomrule
\end{longtable}

这条曲线的形状与第十章的 roofline 是同一族：拐点左侧被深度（在途
数）限制，右侧被设备能力封顶，拐点本身由 $\lambda W$ 决定。工程读
法也只有一句：深度按 $\lambda W$ 取，再留一档余量；过了拐点一分钱
吞吐也买不到，买到的全是时延。

主机侧的经验值：每队列深度取 32 至 128 是常见区间；\code{fio} 压测
的惯例 \code{iodepth=32} 正在典型拐点附近，这也是它成为默认值的
原因。时延敏感的负载反而要浅：日志刷盘、同步写这类请求应走独立队列
或浅队列，避免排在批量流量后面——用第十章的语言说，浅队列买的是时
延的可预期性，不是平均速度。多队列改变的是这笔账的分母：$n$ 个核各
开一对队列，每对只承担 $\lambda/n$，同样的总吞吐用更浅的每队深度就
够，排队时延同步下降——NVMe 的多队列与深度取舍，是同一条定律的两
个应用。

\topic{调试案例两则}
两则案例的共同点是：症状都出现在吞吐或时延上，根因都在队列机制与介
质特性的交界处，定位方法都是"先看队列计数，再看介质状态"。

案例一：队列深度 1 的吞吐瓶颈。症状：新服务器的 NVMe SSD 规格
3.5 GB/s，实测顺序读只有约 40 MB/s，4K 随机读约 1 万 IOPS，离规格
差两个数量级。定位：\code{iostat -x} 显示设备利用率 100\%，但队列
长度 \code{aqu-sz} 恒为 1.00；压测用的是 \code{fio} 默认的同步
I/O，\code{iodepth=1}。根因：同步 I/O 一次只有一个请求在途，
Little 定律直接封顶——$4\ \text{KB}/100\ \mu\text{s}=40$ MB/s。设
备没有病，是提交方式只给它一单做一单。修复：改用 \code{libaio} 或
\code{io\_uring}，\code{iodepth} 提到 32。同机复测：4K 随机读从 1
万 IOPS 升到约 32 万 IOPS，顺序读回到 3 GB/s 量级——吞吐涨了约 32
倍，硬件一分钱没动。教训：报外设性能必须先报队列深度，没有深度的
IOPS 数字没有意义；看到利用率 100\% 先看队列长度——利用率满而队列
空，瓶颈在提交侧，不在设备侧。

案例二：GC 风暴导致的周期性延迟尖峰。症状：持续随机写的服务，P50
写时延稳定在 200 $\mu$s，但每隔几秒出现一批 10 至 50 ms 的慢写，
P99 从 1 ms 恶化到 30 ms；写入速率越快，尖峰间隔越短。定位：时延直
方图呈双峰；盘已写到接近满容量；SSD 健康信息里主机写入量远小于
NAND 实际写入量，写放大超过 3。根因：盘写满后没有现成空闲块，每次
写都要先做一次 GC——读出整块、搬走有效页、擦除——写请求排在一次
毫秒级的擦除后面；空闲块耗尽一批就尖峰一次，所以呈周期性，且写入越
快耗尽越快。对照实验：对盘做安全擦除（或全量 TRIM）后重测，尖峰消
失、P99 回落，确认根因在空闲块供给，不在队列深度或接口。修复：文件
系统开启 \code{discard} 或定期 \code{fstrim}，让删除的文件及时变回
设备可用的空闲块；预留空间提到 20\% 以上；时延敏感的服务不把盘写
满，容量水位线设在 70\%。用闲置容量换空闲块的持续供给，本质上是给
GC 留余量——与第十章"服务台保持余量"是同一条纪律。

两则案例放在一起看，方法上的共性比结论本身更值得带走。两案的设备都
是完好的，换硬件解决不了任何一案；两案的突破口都是计数器而不是猜
测——队列长度、利用率、写放大、时延直方图，全是现成的可观测数据；
两案的修复也都不动设备，一个改提交方式，一个改介质水位。外设问题的
默认值应当是"机制问题"而不是"硬件坏了"：先把 Little 定律和介质的
脾气代进去算一遍，算不通再怀疑硬件。

\topic{外设问题的通用读法}
全章收束为一张检查表。面对任何外设的性能或异常问题，按三步走：先问
介质特性——这类设备的物理限制是什么；再问队列机制——它用什么结构
弥补介质的短板；最后测量验证——不看规格书看计数器，不看平均看百分
位。三类设备放进同一张表：

\begin{longtable}{L{0.10\linewidth}L{0.26\linewidth}L{0.26\linewidth}L{0.26\linewidth}}
\toprule
设备 & 介质先问 & 队列再问 & 测量验证 \\
\midrule
SSD & 空闲块剩多少，写放大多少 & 提交深度多少，谁在排队 & \code{fio} 按百分位报时延，\code{iostat} 看 \code{aqu-sz} \\
机械盘 & 随机还是顺序，寻道占几成 & 调度器合并了吗，NCQ 用上了吗 & \code{iostat} 看 await 与服务时间之差 \\
网卡 & 先到 pps 顶还是带宽顶，帧多大 & 聚合参数多少，多队列吃满核了吗 & \code{ethtool -S} 看丢包计数，按核看中断分布 \\
\bottomrule
\end{longtable}

三类设备共享同一张检查表不是巧合：第五章的描述符环与 DMA、第十章的
Little 定律与百分位纪律，在每一块外设上都是同一套语言。读外设问题
的功力，说到底就是先认出介质那一栏，再走完后两栏。

把三步走完整演练一次。某数据库机器报告"磁盘慢"：第一步问介质——这
是一块 SATA SSD 还是 NVMe SSD、写到了几成满、随机写占多大比例，因
为满盘加随机写正是 GC 尖峰的经典触发条件；第二步问队列——提交深度
多少、走的是单队列还是多队列、同步写有没有和批量读混在同一对队列
里，因为平均时延 200 $\mu$s 但 P99 到 20 ms 的形状，指向的是"排在
慢事件后面"而不是"介质本身慢"；第三步测量验证——\code{fio} 单独复
现负载并按百分位报时延，\code{iostat -x} 对照 \code{aqu-sz} 与利用
率，再查 SSD 健康信息里的写放大与剩余备用块。三步走完，"磁盘慢"这
句模糊报告就被翻译成一句可执行的话：盘已到 92\% 容量、写放大 4.5，
先把水位降到 70\% 再谈别的。同一个流程搬到网卡上，只是第一栏换成
pps 与帧长、第三栏换成 \code{ethtool} 计数——方法不动，参数换设备。

\section*{章末练习}

\begin{longtable}{L{0.2\linewidth}L{0.5\linewidth}L{0.2\linewidth}}
\toprule
任务 & 要求 & 学习目标 \\
\midrule
扇区对齐 & 某 4K 扇区硬盘的分区从 63 号逻辑（512 B）扇区开始，文件系统按 4 KB 块写入。求分区起始字节偏移是否为 4 KB 的整数倍；说明每次文件系统块写要触及几个物理扇区、代价是什么，并给出修复做法 & 分区边界与物理扇区对齐 \\
写放大计算 & 某 SSD 页 16 KB、每块 256 页，随机写稳态下每块平均 75\% 为有效页。求写放大系数；主机累计写入 10 GB 时，NAND 实际被写入多少 & 写放大 $=1/(1-\text{有效占比})$ 的来源 \\
随机 IOPS 估算 & 7200 rpm 硬盘，平均寻道 4 ms，顺序传输率 150 MB/s。估算 4 KB 随机读的 IOPS 与等效带宽，并与顺序读带宽比较 & 毫秒级寻道加旋转决定随机 IOPS \\
中断聚合权衡 & 1 Gb/s 网卡满载 64 B 小帧（线速约 148.8 万 pps），每次中断处理约 5 $\mu$s。求不聚合时的 CPU 开销；聚合到每 32 帧一次中断后的开销；说明聚合付出的代价 & 聚合用时延换 CPU 占用 \\
NVMe trace 排序 & 把一次 NVMe 读的这些事件排成时间序：DMA 写数据到内存、写 SQ 尾 doorbell、页缓存未命中、设备写 CQ 完成项、MSI 中断、驱动写命令项到 SQ、唤醒进程 & 提交—执行—完成三段的次序 \\
队列深度选择 & 某 NVMe SSD 4 KB 随机读时延 90 $\mu$s，应用目标 40 万 IOPS。用 Little 定律求所需最小队列深度；若设备服务率上限 50 万 IOPS，求 QD=256 时的平均时延 & $L=\lambda W$；过拐点只加排队 \\
\bottomrule
\end{longtable}

\section*{参考解答与要点}

\begin{longtable}{L{0.18\linewidth}L{0.5\linewidth}L{0.22\linewidth}}
\toprule
任务 & 参考解答 & 必须掌握的要点 \\
\midrule
扇区对齐 & 起始偏移 $=63\times512=32256$ B，$32256/4096=7.875$，不是 4 KB 的整数倍，分区未对齐。文件系统每个 4 KB 块都跨越两个物理 4 KB 扇区，写一块要先读出两个扇区、修改、再写回（读-改-写），写代价约翻倍；SSD 上同理跨页、放大 GC 负担。修复：分区起始对齐到 1 MiB（2048 个 512 B 扇区）边界，现代分区工具默认如此 & 先算偏移再谈对齐；不对齐的代价是读-改-写 \\
写放大计算 & 每回收一个块：块容量 $256\times16$ KB $=4$ MB，其中有效 $4\times0.75=3$ MB 须先搬出，只换得 1 MB 空闲。稳态下主机每写 1 MB，NAND 要写 1 MB 数据加 3 MB 搬运共 4 MB，写放大 $=4=1/(1-0.75)$。主机累计写 10 GB，NAND 实际写入 $10\times4=40$ GB & 有效占比越高写放大越大；留空闲空间直接降放大 \\
随机 IOPS 估算 & 7200 rpm 每转 $60/7200\approx8.33$ ms，平均旋转等待半转 $4.17$ ms；加平均寻道 4 ms，4 KB 传输 $4096/(1.5\times10^{8})\approx0.03$ ms，单次共约 8.2 ms。IOPS $\approx1/0.0082\approx122$；等效带宽 $122\times4$ KB $\approx0.49$ MB/s，与顺序 150 MB/s 相差约 300 倍 & 随机 IOPS 由寻道加旋转写死；顺序与随机差两个数量级 \\
中断聚合权衡 & 不聚合：$1.488\times10^{6}\times5\ \mu\text{s}\approx7.4$ CPU 秒/秒，约需 8 个核专职处理中断，系统实际早已丢包。聚合 32 帧一次：中断率 $1.488\times10^{6}/32\approx4.65\times10^{4}$ 次/秒，CPU 开销 $\approx4.65\times10^{4}\times5\ \mu\text{s}\approx0.23$ 秒/秒，约一个核的 23\%。代价：帧要等够一批或等到聚合定时器超时，平均时延增加约半个聚合窗口（定时 64 $\mu$s 时约 32 $\mu$s） & 聚合把到达率除以批大小；吞吐与时延的交换 \\
NVMe trace 排序 & 次序：页缓存未命中 → 驱动写命令项到 SQ → 写 SQ 尾 doorbell → 设备 DMA 写数据到内存 → 设备写 CQ 完成项 → MSI 中断 → 唤醒进程。两处次序是硬性约定：doorbell 必须在命令项写好之后（先可见、再通知），CQ 完成项必须在 DMA 数据之后（先数据、后完成），唤醒只能在中断处理确认完成之后 & 提交—执行—完成三段；每段先后都是可见性约定 \\
队列深度选择 & $L=\lambda W=4\times10^{5}\times90\times10^{-6}=36$，最小深度 36，工程上取 64 留一档余量。设备上限 50 万 IOPS 时已经饱和，QD=256 的平均时延 $\approx256/(5\times10^{5})=512\ \mu$s：吞吐仍是 50 万 IOPS，多出的深度全部变成排队，时延约为 90 $\mu$s 的 5.7 倍 & 深度按 $\lambda W$ 取并留余量；过拐点后加深只买排队 \\
\bottomrule
\end{longtable}
